\documentclass[11pt,a4paper]{article}
\usepackage[T1]{fontenc}
\usepackage[utf8]{inputenc}
\usepackage{lmodern,amsmath,amssymb,amsthm,mathtools}
\usepackage[margin=25mm]{geometry}
\usepackage{booktabs,longtable,array,enumitem,microtype}
\usepackage[hidelinks]{hyperref}
\setlength{\parindent}{0pt}
\setlength{\parskip}{5pt plus 1pt}
\setlist{leftmargin=1.7em,itemsep=3pt,topsep=4pt}
\newtheorem{theorem}{Theorem}[section]
\newtheorem{proposition}[theorem]{Proposition}
\newtheorem{lemma}[theorem]{Lemma}
\newtheorem{corollary}[theorem]{Corollary}
\theoremstyle{definition}
\newtheorem{definition}[theorem]{Definition}
\newtheorem{example}[theorem]{Example}
\newtheorem{remark}[theorem]{Remark}
\newcommand{\Set}{\mathbf{Set}}
\newcommand{\Cat}{\mathbf{Cat}}
\newcommand{\Biprof}{\mathbf{Biprof}}
\newcommand{\V}{\mathcal V}
\newcommand{\Lcat}{\mathcal L}
\newcommand{\K}{\mathcal K}
\newcommand{\J}{\mathcal J}
\newcommand{\D}{\mathcal D_{\V}}
\newcommand{\A}{\mathcal A}
\newcommand{\Uclass}{\mathcal U}
\newcommand{\Wclass}{\mathcal W}
\newcommand{\Aug}{\operatorname{AugFam}}
\newcommand{\id}{\mathrm{id}}
\newcommand{\List}{\operatorname{List}}
\newcommand{\im}{\operatorname{im}}
\newcommand{\End}{\operatorname{End}}
\newcommand{\Rel}{\operatorname{Rel}}
\newcommand{\FinVect}{\mathbf{FinVect}}
\newcommand{\FinSet}{\mathbf{FinSet}}
\newcommand{\Ab}{\mathbf{Ab}}
\newcommand{\Grp}{\mathbf{Grp}}
\newcommand{\psint}{\int_{\mathrm{ps}}}
\newcommand{\inl}{\mathrm{inl}}
\newcommand{\inr}{\mathrm{inr}}
\newcommand{\starc}{\mathbin{\star}}
\newcommand{\diamondc}{\mathbin{\diamond}}
\title{General Independent Interaction Criterion\\[5pt]
\large Augmented rig contexts and relative occurrence displayed content}
\date{1 November 2026}
\begin{document}
\maketitle

\begin{abstract}
An independently given coherent symmetric rig category and an independently
valid contentad determine a complete retained-cell interaction when the content
has a relative polynomial occurrence display, cartesian unit and accumulation,
and Beck--Chevalley transport. The context doctrine, its noninvertible
coefficient maps, dependent accumulation and finite occurrence tensor are
constructed from the rig; no mixed lifting is supplied. The content input is
not restricted to a receipt recipe. The proof derives the crossing, all native
cell actions and pure endpoint squares, pure-middle balance, four invertible
modifications and all ten pseudodistributive-law diagrams. It does not assume
that set carriers detect coefficient maps, or that content extension units and
multiplications are invertible.

Marked exhaustive occurrence normal forms force display exactness, cartesian
unit and accumulation, and canonical Beck--Chevalley along context-realizable
projections. This gives an exact relative converse when those projections are
rich enough, and an explicit necessary/sufficient gap otherwise. Strong input
changes and exact one-sided controls are proved. Distinct content mechanisms,
linear resources and finite probabilistic and consumable-cost protocols test
the same criterion. A separate abelian display/codisplay theorem and a sharp
bounded group obstruction record the earlier pullback/pushout branch.
The result is offered for O1 acceptance at the stated structural scope, not as
an arbitrary-pair classification or completion of the expanded contract.
\end{abstract}

\tableofcontents

\section{Scope and principal result}\label{sec:scope}

The earlier \emph{Indexed Receipt Interaction Theorem}~\cite{Receipt} proved a
scoped family: normalized semiring-weighted histories against content generated
from a polynomial monad by indexed receipts. The present result continues the
broader O1 requirement~\cite{Contract}. Both inputs now vary independently:
resources range over coherent symmetric rig categories, and content ranges over
already valid contentads satisfying stated content-only exactness conditions.
Receipts are one example, not the definition of admissible content. The earlier
proof remains valid and is recovered, rather than overwritten.

The theorem is deliberately a family criterion. It does not say that arbitrary
contextads and contentads interact, classify all possible interactions, or find
globally weakest assumptions. The bounded converse fixes explicit witness and
projection marks; it does not assume the desired mixed comparison as a mark.
The remaining restrictions are collected in Section~\ref{sec:acceptance}.

\begin{theorem}[Independent relative interaction]\label{thm:main}
Choose the following data before constructing any mixed composition.
\begin{enumerate}[label=\textup{(\roman*)}]
\item A class $\J$ of set maps containing identities and isomorphisms and
closed under pullback and composition; a finitary cartesian polynomial monad
$T$ with $T(\J)\subseteq\J$.
\item A locally small coherent symmetric rig category $\V$. Form augmented
finite $\V$-families and restrict indexed context grades to those whose carrier
projection belongs to $\J$.
\item An independently valid contentad $(p:\Lcat\to\Set,d,U,a,\eta^L,\mu^L)$,
with its actual global opcartesian transport, genuine cartesian lifts over
$\J$, and a natural occurrence map $\nu:d\Rightarrow Tp$ satisfying the unit and
accumulation equations. Require display pullbacks, cartesian preservation by
$U$ and $a$, and invertible canonical Beck--Chevalley mates precisely as in
Definition~\ref{def:content}.
\end{enumerate}
These data determine a full pseudodistributive law
\[
 \tau:N\odot M\Longrightarrow M\odot N
\]
on the original category-valued modules, retaining all native noninvertible
cells and arbitrary pure endpoint squares. It induces a relative
pseudopromonad on $H(A,B)=(e_A\downarrow d_B)$.

Within the exhaustive addressed occurrence normal forms of
Definition~\ref{def:marked}, an interaction forces the corresponding display,
cartesian preservation and Beck--Chevalley conditions along the realizable
context projections $\D$. If independently $\D=\J$, these conditions are
necessary and sufficient in that marked class. Otherwise the necessary and
sufficient classes remain distinguished.

The construction is strong under the parameter maps of
Definition~\ref{def:parameters}. The pure-content inclusion is strongly
locally fully faithful. The pure-context inclusion is strong and is locally
fully faithful at all interfaces if and only if
$\End_{\Lcat_B}(U_B)=\{\id\}$ for every $B$.
\end{theorem}

Proofs appear in Sections~\ref{sec:rig}--\ref{sec:controls}. The independently
verified examples are in Section~\ref{sec:instances}. This is a mathematical
proof dossier, not proof-assistant certification. The supplied source is
intentionally uncompiled; finite regressions and source-integrity checks have
the limited roles described in Appendix~\ref{app:checks}.

\section{Ambient categories and an assumption ledger}\label{sec:ambient}

\subsection{Size and retained relative composition}
Fix nested universes. Sets and all individual fibres lie in the first suitable
universe; categories of such sets, category-valued modules, and their relative
pseudocoends are formed at successively larger levels. No category of all sets
is asserted to be small. Every witness carrier and every polynomial position
set below is finite. Interface sets, polynomial shape sets and content-grade
carriers need not be finite.

A module is a pseudofunctor $X:\Set^{\mathrm{op}}\times\Set\to\Cat$.
Write its relative tensor in execution order:
\[
 (X\odot Y)(A,C)=\psint^{B\in\Set}X(A,B)\times Y(B,C).
\]
The concrete pseudocoend has objects $[B;x,y]$, arrows from the product
categories and invertible pure-middle balance arrows, with naturality, identity
and composition relations. Native arrows are not inverted. Iterated tensors
use the common flattened, multiply balanced presentation. This gives the
ambient associators, unitors and interchangers; one must check a proposed law
on those generators, not only on fixed-middle objects.

For indexed contexts $(\K,e)$ and content $(\Lcat,d)$ the modules are
\[
 M(A,B)=(e_A\downarrow B),\qquad
 N(A,B)=(A\downarrow d_B),\qquad
 H(A,B)=(e_A\downarrow d_B).
\]
An object of $H(A,B)$ is $(P,Q,f:e_AP\to dQ)$, and a fixed-endpoint cell is
\begin{equation}\label{eq:comma}
 (\alpha:P\to P',\ \beta:Q\to Q'),\qquad
 d\beta\,f=f'e_A\alpha.
\end{equation}
For endpoint maps $u:A\to A'$ and $v:B\to B'$, a square has
\begin{equation}\label{eq:endpoints}
 \alpha:P\to u^*P',\quad \beta:v_!Q\to Q',\quad
 d\beta\,d\kappa_{v,Q}\,f
 =f'e_u(P')e_A\alpha,
\end{equation}
where $\kappa_{v,Q}:Q\to v_!Q$ is the chosen opcartesian arrow and
$e_u(P'):e_Au^*P'\to e_{A'}P'$ is the context extension map. These are
actual grade morphisms, including parallel and noninvertible ones.

\begin{lemma}[Retained comma normalization]\label{lem:normalization}
There is a module equivalence $M\odot N\simeq H$ which retains all native
cells and endpoint squares.
\end{lemma}
\begin{proof}
Send $[B;(P,f:eP\to B),(Q,g:B\to dQ)]$ to $(P,Q,gf)$.
The opposite functor uses the representative
$[eP;(P,\id_{eP}),(Q,f)]$ for $(P,Q,f)$. For a cell $(\alpha,\beta)$ satisfying~\eqref{eq:comma}, first apply the
native $N$-arrow $\beta$ to obtain
$[eP;(P,\id),(Q',f'e\alpha)]$. Next use the balance arrow along
$e\alpha:eP\to eP'$ to obtain
$[eP';(P,e\alpha),(Q',f')]$. Finally apply the native $M$-arrow
$\alpha:(P,e\alpha)\to(P',\id)$. Thus both native components, as well as
the change of middle object, are explicitly retained. Naturality
and composition are the balance relations. The composite to $H$ is the
identity, and the other composite is naturally isomorphic to the identity by
the balance arrow for $f:eP\to B$. Its naturality and the two triangle
identities are again the defining balance relations. Pure endpoint actions
commute with these constructions by their own pseudofunctor constraints.
Thus no quotient by isomorphism classes, support, or connected components is
used. This is the inherited normalization~\cite{Unified}, spelled out here to
fix what the new law must preserve.
\end{proof}

The one-sided context and content presentations determine relative
pseudopromonads $M,N$. Once the full law is proved, established
pseudomonad-composition machinery~\cite{Marmolejo,GambinoLobbia} gives
\begin{equation}\label{eq:composite}
 j_H=(j_M\odot j_N),\qquad
 m_H=(m_M\odot m_N)(M\odot\tau\odot N),
\end{equation}
transported through Lemma~\ref{lem:normalization}. This inherited step is not
a new compatibility assumption: the work is to construct and verify $\tau$.

\subsection{Ledger of material assumptions}
The roles below distinguish independent inputs from derived conclusions.
All chosen lift comparisons are the comparisons forced by their universal
properties, not freely supplied mixed coherence.

\begin{longtable}{@{}p{.08\linewidth}p{.37\linewidth}p{.27\linewidth}p{.20\linewidth}@{}}
\toprule
ID&Statement and scope&Role and proof uses&Status and boundary\\
\midrule\endhead
S1&Nested universes; finite witnesses and polynomial positions.&Typing and
existence of indexed sums and tensors.&Supplied. Infinite occurrences need
additional hypotheses.\\
J1&$\J$ contains identities and isomorphisms, is pullback/composition stable;
$T\J\subseteq\J$.&Closure of context grades and every relative exactness use.&
Supplied independently. Realizable maps may be fewer.\\
V1&Two symmetric monoidal structures with coherent invertible distributors
and absorption.&Constructs all coefficient maps, dependent accumulation and
occurrence tensor.&Supplied rig axioms; no coproduct property assumed.\\
L1&An already valid one-sided contentad, with actual covariant transport and
possibly noninvertible $\eta^L,\mu^L$.&One-sided pseudo-unit and associativity
laws; underlying $N$.&Supplied and checked separately. The slice false
example fails here.\\
L2&Global opfibration with the given content transport, and genuine cartesian
lifts over $\J$.&All native context-map factorizations, including arbitrary
maps into a $\J$-lift.&Supplied. No cartesian lift outside $\J$ or
opcartesian preservation by $d$ assumed.\\
L3&Natural $\nu:d\Rightarrow Tp$ and its unit/accumulation equations.&Shapes,
labels and addressed substitution.&Supplied content-only condition. Strict
arity-changing cells excluded.\\
L4&Cartesian display over $\J$.&Crossing body and body uniqueness.&Supplied;
forced on $\D$ by marked exhaustive recognition.\\
L5&$U$ and typed accumulation preserve cartesian arrows over $\J$.&Derives
invertible native grade unit/multiplication comparisons.&Supplied
content-only property; marked recognition forces it on $\D$.\\
L6&Canonical BC invertible for pullbacks with vertical maps in $\J$.&Pure-middle
balance, arbitrary horizontal pure maps.&Supplied content-only property;
marks force the canonical mate.\\
C1&$\End_{\Lcat_B}(U_B)$ trivial.&Exactly the extra pure-context embedding
control.&Optional; necessary and sufficient globally.\\
R1&Exhaustive invertible crossing body and original grade-projection marks.&
Defines the bounded recognition class.&Not an assumption of the sufficient
construction; not inferred for arbitrary laws.\\
P1&Strong rig/contentad maps, cartesian monad map and extension comparison.&
Parameter functoriality.&Separate input maps; mixed comparison derived.\\
\bottomrule
\end{longtable}

\section{Augmented rig families and their full native calculus}\label{sec:rig}

\subsection{Structural coherence with addressed occurrences}
A coherent symmetric rig category has symmetric monoidal structures
$(\oplus,0)$ and $(\otimes,I)$, invertible left/right distributors and
absorption isomorphisms satisfying the distributivity coherence axioms.
Neither $\oplus$ nor $0$ is assumed to have a categorical coproduct or initial
universal property. Choose finite indexed sums and tensor products, with
specified comparison maps for partitions, permutations and rebracketing.

Here and below, a structural map is identified by its permutation of
\emph{tagged occurrences}, not merely by unlabelled endpoints. The facts used
are partition associativity, tensor coherence, finite distributivity and its
compatibility with partitions, grouping, permutations, zeros and units.
For clarity about the source, Laplaza's Proposition~10~\cite{Laplaza} applies
to regular formal expressions, not arbitrary diagrams sharing endpoints.
In our expressions each branch and independent occurrence first receives a
separate formal variable. Expanded nonzero monomials are squarefree and
distinct, so the required regular coherence applies. Distinct formal variables
may subsequently be interpreted by the same actual object.

To include inverse distributors, fix each expression's forward expansion to
its chosen sum-of-monomials normal form. Regular coherence equates the two
forward paths across a distributor. Since the distributor is invertible, the
same equation holds across its inverse. Induction on a path, followed by
cancellation of the target normalization isomorphism, compares any two paths
with the specified permutation. Unit-only and zero cases use unit and
absorption coherence. For diagrams containing arbitrary arrows, first use
naturality and bifunctoriality to move those arrows through the structural
maps; apply structural coherence only to the remaining structural paths.
In particular a swap of two equal actual summands is not automatically the
identity.

\subsection{The category of augmented families}
\begin{definition}
An augmented family is $P=(X,(V_x)_{x\in X},\epsilon_P)$, where $X$ is finite
and $\epsilon_P:S(P):=\bigoplus_xV_x\to I$ is one global augmentation.
A native arrow $h:P\to P'$ consists of $f:X\to X'$ and actual maps
\[
 \alpha_{x'}:\bigoplus_{fx=x'}V_x\longrightarrow V'_{x'}
\]
satisfying $\epsilon_{P'}A_h=\epsilon_P$. Here $A_h:S(P)\to S(P')$
partitions by $f$ and then applies $\bigoplus_{x'}\alpha_{x'}$.
\end{definition}
No component maps $V_x\to I$ are assumed: they cannot generally be extracted
from a global augmentation without coproduct injections.

For $k=(g,\beta):P'\to P''$, define the component of $kh$ at $z$ by
\[
 \bigoplus_{gfx=z}V_x
 \cong\bigoplus_{gx'=z}\bigoplus_{fx=x'}V_x
 \xrightarrow{\oplus\alpha_{x'}}\bigoplus_{gx'=z}V'_{x'}
 \xrightarrow{\beta_z}V''_z.
\]
Then $A_{kh}=A_kA_h$ by partition coherence, so augmentation is preserved.
Both bracketings of a triple composite give the same nested sums of the
actual maps followed by their compositions. Additive coherence and
bifunctoriality prove associativity; singleton comparison maps give identities.
Thus these objects and arrows form $\Aug(\V)$.

\begin{lemma}\label{lem:terminal}
$J_\V=(\{*\},I,\id_I)$ is terminal in $\Aug(\V)$.
\end{lemma}
\begin{proof}
The carrier map is unique, and the single resource component is forced by
the augmentation equation to be $\epsilon_P$.
\end{proof}
Empty families exist exactly when a map $0\to I$ is available. The carrier
functor $\Aug(\V)\to\FinSet$ need not be faithful.

For a set $A$, put $\K_A=\Aug(\V)^A$ and
\[
 e_AP=\coprod_{a\in A}|P_a|,\qquad \pi_P(a,x)=a.
\]
Reindexing is pointwise substitution, so extension reindexing squares are
pullbacks. Restrict to $\pi_P\in\J$. All native family arrows between such
objects remain admitted; their carrier maps themselves need not be in $\J$.
The context unit is the constant family $J_\V$.

\subsection{Dependent accumulation including aggregation maps}
For $P_a=(X_a,V_{ax},\epsilon_a)$ and
$Q_{(a,x)}=(Y_{ax},W_{axy},\delta_{ax})$, define
\[
 |(P\starc Q)_a|=\{(x,y):x\in X_a,\ y\in Y_{ax}\},\qquad
 (P\starc Q)_{a,(x,y)}=V_{ax}\otimes W_{axy}.
\]
Its augmentation is
\begin{equation}\label{eq:depaug}
 \bigoplus_{x,y}V_{ax}\otimes W_{axy}
 \cong\bigoplus_x\left(V_{ax}\otimes\bigoplus_yW_{axy}\right)
 \xrightarrow{\oplus(\id\otimes\delta_{ax})}
 \bigoplus_x(V_{ax}\otimes I)
 \cong\bigoplus_xV_{ax}\xrightarrow{\epsilon_a}I.
\end{equation}
The extension is the dependent-sum reassociation
$e_A(P\starc Q)\cong e_{e_AP}Q$. Its projection is a composite of the two
input projections, so belongs to $\J$.

Suppress $a$. An arrow in the category of dependent pairs consists of
$(f,\alpha):P\to P'$ and
$(g_x,\beta_x):Q_x\to Q'_{fx}$. Its accumulated carrier map is
$(x,y)\mapsto(fx,g_xy)$. At $(x',y')$ its resource map is
\begin{align}\label{eq:depmap}
 &\bigoplus_{fx=x'}\bigoplus_{g_xy=y'}V_x\otimes W_{xy}
 \cong\bigoplus_{fx=x'}\left(V_x\otimes
                  \bigoplus_{g_xy=y'}W_{xy}\right)\notag\\
 &\xrightarrow{\oplus(\id\otimes\beta_{x,y'})}
 \bigoplus_{fx=x'}(V_x\otimes W'_{x'y'})
 \cong\left(\bigoplus_{fx=x'}V_x\right)\otimes W'_{x'y'}
 \xrightarrow{\alpha_{x'}\otimes\id}V'_{x'}\otimes W'_{x'y'}.
\end{align}
The $x$-dependent $\beta_x$ are applied \emph{before} $\alpha$ merges their
outer fibre. This ordering is essential.

\begin{lemma}\label{lem:depfunctor}
Dependent accumulation is functorial on all native dependent-pair arrows and
preserves their augmentations.
\end{lemma}
\begin{proof}
Sum~\eqref{eq:depmap} over $(x',y')$ and move structural maps by naturality.
The composite with the target augmentation becomes
\[
 \bigoplus_x[V_x\otimes S(Q_x)]
 \xrightarrow{\oplus(\id\otimes A_{k_x})}
 \bigoplus_x[V_x\otimes S(Q'_{fx})]
 \xrightarrow{\oplus(\id\otimes\delta'_{fx})}
 \bigoplus_xV_x\xrightarrow{A_h}\bigoplus_{x'}V'_{x'}
 \xrightarrow{\epsilon'}I.
\]
Use $\delta'_{fx}A_{k_x}=\delta_x$ and $\epsilon'A_h=\epsilon$ to recover
\eqref{eq:depaug}.

For composition, introduce a second outer map $(f',\alpha')$ and
continuations $(g'_{x'},\beta'_{x'})$. At a fixed final fibre, distribute all
intermediate sums. The essential operations are the sums of
$\id_{V_x}\otimes\beta_x$, then
$\id_{V_x}\otimes\beta'_{fx}$, then $\alpha_{fx}\otimes\id$, then
$\alpha'\otimes\id$. The route accumulating each arrow separately places
$\alpha_{fx}$ before $\beta'_{fx}$. Those two operations commute by tensor
bifunctoriality because $\beta'_{fx}$ is constant on the fibre merged by
$\alpha_{fx}$. All other changes are distributivity naturality and partition
coherence. Thus both routes have exactly the same actual resource composite.
Empty preimages are included: absorption naturality identifies $0\otimes b$
with $\id_0$ after absorption, so the same interchange applies even to
nonunique maps $0\to V'$. No initiality argument is used. Identities follow
from the singleton comparisons and functoriality of tensor.
\end{proof}

The associator has carrier $((x,y),z)\mapsto(x,(y,z))$ and resource map the
tensor associator. Distributing by $x,y,z$, both augmentation composites first
apply the augmentation of the innermost continuation, then the middle, then
the outer augmentation. Naturality and monoidal coherence identify them.
The same normalization with~\eqref{eq:depmap} proves naturality. Unit arrows
remove a singleton choice and apply tensor unitors. Pentagon and triangle
compare identical tuple maps and tensor structural maps, and hence commute.
Pure base change is pointwise substitution. These are the full contextad
operations, not merely operations on objects or on a core.

\subsection{Independent occurrence tensor and dependent Fubini}
For a finite addressed set $E$ and augmented families $P_i$, define
\begin{equation}\label{eq:tensor}
 \left|\boxtimes_{i\in E}P_i\right|=\prod_i|P_i|,\qquad
 \left(\boxtimes_iP_i\right)_x=\bigotimes_iV_{i,x_i}.
\end{equation}
Its augmentation distributes to $\bigotimes_iS(P_i)$ and applies
$\bigotimes_i\epsilon_i$, followed by $\bigotimes_iI\cong I$.
For $E=\varnothing$ it is $J_\V$. A family of arrows $(f_i,\alpha_i)$ acts
at a target tuple $y$ by
\begin{equation}\label{eq:tensormap}
 \bigoplus_{f_ix_i=y_i\ \forall i}\bigotimes_iV_{i,x_i}
 \cong\bigotimes_i\left(\bigoplus_{f_ix_i=y_i}V_{i,x_i}\right)
 \xrightarrow{\bigotimes_i\alpha_{i,y_i}}\bigotimes_iV'_{i,y_i}.
\end{equation}
The tensor of the input augmentation equations proves augmentation preservation.
Distributing intermediate fibre sums and using tensor bifunctoriality proves
functoriality. Permutation and regrouping maps are natural by distributor and
tensor naturality and coherent by the tagged structural argument above.
This operation is not asserted to be a categorical product, and needs no
resource diagonal $V\to V\otimes V$.

\begin{proposition}[Full native dependent Fubini]\label{prop:fubini}
There is a natural augmented isomorphism
\[
 \Phi:\boxtimes_i(P_i\starc Q_i)
 \xrightarrow{\sim}
 (\boxtimes_iP_i)\starc\bigl(x\mapsto\boxtimes_iQ_{i,x_i}\bigr).
\]
Its carrier is $((x_i,y_i))_i\mapsto((x_i)_i,(y_i)_i)$ and its coefficient
map is the addressed shuffle
\[
 \bigotimes_i(V_{i,x_i}\otimes W_{i,x_i,y_i})
 \cong(\bigotimes_iV_{i,x_i})\otimes(\bigotimes_iW_{i,x_i,y_i}).
\]
It is natural in every aggregated native map, and compatible with dependent
associativity, independent regrouping, units and permutations.
\end{proposition}
\begin{proof}
After distribution, both augmentations first apply all continuation
augmentations $\delta_{i,x_i}$ and then all outer augmentations $\epsilon_i$.
The only interchanges are on disjoint tensor factors. Naturality of
distribution and tensor symmetry proves equality of these actual maps.

For arbitrary outer maps $(f_i,\alpha_i)$ and continuation maps
$(g_{i,x},\beta_{i,x})$, fix a target tuple $(x'_i,y'_i)$. Both routes start
from the sum over $f_ix_i=x'_i$ and $g_{i,x_i}y_i=y'_i$. Distribute by $i$,
apply each $\beta_{i,x_i}$ before merging by $\alpha_{i,x'_i}$, and regroup.
Both routes reduce to the tensor over $i$ of the actual composite
in~\eqref{eq:depmap}. Distributor naturality and bifunctoriality identify
these reductions, and the final tensor shuffle has the same tagged
permutation. This proves equality of coefficient arrows, not just carriers.

For all stated coherences expand to branch-indexed tensor monomials. A monomial
records each ancestor choice and a distinct tag for each occurrence. It is
squarefree; distinct histories give distinct formal monomials. Each structural
route produces the same target nesting and permutation, so regular coherence
applies. With native arrows, move them through structure by naturality first.
The remaining exchanges act on disjoint factors or are exactly the
constant-on-an-aggregated-fibre exchange in Lemma~\ref{lem:depfunctor}.
Augmentations remain in continuation-before-parent order. Empty occurrence
sets use unit coherence and empty branch sets use absorption coherence.
\end{proof}

\section{Independent relative content exactness}\label{sec:content}

\subsection{Typed input and its one-sided meaning}
Fix a presentation
\[
 TB=\coprod_{s\in S}B^{E_s},\qquad |E_s|<\infty.
\]
A cartesian polynomial monad has cartesian natural transformations
$\eta^T:\id\Rightarrow T$ and $\mu^T:TT\Rightarrow T$.
Cartesianness of $\eta^T$ says its selected shape has one position.
Cartesianness of $\mu^T$ gives a bijection
\begin{equation}\label{eq:leaves}
 E_{m(s,(s_i))}\xrightarrow{\sim}\coprod_{i\in E_s}E_{s_i}.
\end{equation}
The monad laws say that successive flattenings have the same shape and the
same addressed leaf bijection. Indeed a polynomial natural transformation is
determined by its shape map and backwards position maps: evaluate each source
representable on its own position set and its identity labelling. Thus equality
of the monad transformations detects these position maps, not only arities.

\begin{definition}[Relative content input]\label{def:content}
Supply an already valid contentad with total category $p:\Lcat\to\Set$,
extension $d:\Lcat\to\Set$, unit $U_B\in\Lcat_B$, accumulation
$Q\diamondc R\in\Lcat_B$ for $Q\in\Lcat_B$ and $R\in\Lcat_{dQ}$, and bodies
\[
 \eta^L_B:B\to dU_B,\qquad
 \mu^L_{Q,R}:dR\to d(Q\diamondc R).
\]
Validity includes the original native-arrow, covariant pure transport,
associativity, pseudo-unit and naturality laws, so the comma module $N$ is an
independently defined relative pseudopromonad. Neither displayed body is assumed
invertible. Require additionally:
\begin{description}[style=nextline]
\item[L2 relative lifts and global transport] $p$ is a global opfibration,
with its given covariant contentad transport, and has chosen genuine
$p$-cartesian lifts over maps in $\J$. Their cartesian universal property is
against arbitrary base arrows, not merely arrows in $\J$; only existence of
lifts is restricted. No cartesian lift outside $\J$ is required.
\item[L3 arity] $\nu:d\Rightarrow Tp$ is natural on all total-category arrows,
and
\begin{equation}\label{eq:arity}
 \nu_{U_B}\eta^L_B=\eta^T_B,\qquad
 \nu_{Q\diamondc R}\mu^L_{Q,R}
 =\mu^T_B\,T\nu_Q\,\nu_R.
\end{equation}
\item[L4 relative display] For a cartesian lift
$c_{f,Q}:f^*Q\to Q$ over $f:X\to B$ in $\J$, the canonical map
\begin{equation}\label{eq:display}
 c^{\mathrm{disp}}_{f,Q}:d(f^*Q)
 \longrightarrow dQ\mathbin{\times_{TB}}TX,
 \qquad z\longmapsto(d c_{f,Q}(z),\nu_{f^*Q}(z))
\end{equation}
is a bijection.
\item[L5 relative cartesian algebra] $U:\Set\to\Lcat$ preserves cartesian
arrows over $\J$. Accumulation preserves cartesian arrows over $\J$ in the
composable-pair category with its relative cartesian lifting structure
\[
 \mathcal E_2=d^*\Lcat\longrightarrow\Lcat\longrightarrow\Set.
\]
An object is $(Q\in\Lcat_B,R\in\Lcat_{dQ})$; an arrow is $(u,v)$ with
$u:Q\to Q'$ over $f$ and $v:R\to R'$ over $du$. Its chosen lift over $f\in\J$ is a cartesian $u$ together with a
cartesian $v$ over $du$; L4 and $T(\J)\subseteq\J$ ensure that this
secondary base map belongs to $\J$.
\item[L6 relative BC] In every pullback square
\[
 \begin{array}{ccc}
 E_u&\xrightarrow{v}&E'\\
 {\scriptstyle\pi_u}\downarrow&&\downarrow{\scriptstyle\pi'}\\
 B&\xrightarrow{u}&B'
 \end{array}
 \quad(\pi_u,\pi'\in\J),
\]
the canonical mate $v_!\pi_u^*\Rightarrow\pi'^*u_!$ is invertible. Horizontal
maps $u,v$ are arbitrary pure maps.
\end{description}
\end{definition}
These conditions refer to $\Lcat,d,T,\J$, not to $\V$, context grades or mixed
composition. Conversely, the rig assumptions mention none of this content.
The content algebra is supplied because it is one of the independent effects;
no law between the two effects is supplied.

\begin{remark}[Why one-sided validity is not optional]\label{rem:slicefalse}
The plausible data $\Lcat_B=\Set/B$, $d(q:Q\to B)=Q$, covariant
postcomposition, unit $\id_B$, and accumulation by composition do not form the
required pseudocontentad with all declared factorization maps. Composing its
purported pure identity with $(b:B\to R,r:R\to C)$ gives the factorization
$(\id_B,rb)$ instead of $(b,r)$. The comparison is $b:B\to R$, which need
not be invertible. For $B=C=1$ and $R=\{0,1\}$ selecting one point, the
middle carriers even have different cardinalities. Equivalently
$f_!(\id_B)=f$ need not be isomorphic to $\id_{B'}$. Pullbacks, a bifibration
and objectwise-looking accumulation do not repair this failed one-sided unit.
Quotienting away the factorization witness would change the required theorem.
\end{remark}

\subsection{Canonical grade comparisons and their coherence}
Put $j_f=d c_{f,Q}$. For $f\in\J$, display identifies $j_f$ with a pullback of
$Tf$, hence $j_f\in\J$. The pair $(f^*Q,j_f^*R)$ is cartesian in
$\mathcal E_2$. L5 and uniqueness of cartesian lifts therefore give
\begin{equation}\label{eq:gradecomparisons}
 U_X\cong f^*U_B,\qquad
 f^*Q\diamondc j_f^*R\cong f^*(Q\diamondc R).
\end{equation}
The directions displayed are the factorizations through the chosen target
cartesian lift; their inverses are equally available. They preserve the
original arrows to $U_B$ and $Q\diamondc R$, not just their arity labels.

Every naturality, identity, composition, unit and associativity diagram for
these comparisons commutes in $\Lcat$: compose both routes with the same
target cartesian arrow. The one-sided contentad laws and naturality identify
the resulting arrows. Cartesian uniqueness then identifies the original
vertical maps. This argument retains parallel native maps even if $d$ is
nonfaithful. Analogously, mate pasting and cartesian/opcartesian factorization
prove coherent composition of the BC comparisons. No preservation of
opcartesian arrows by $d$ is required.

\subsection{Optional detection criterion}
\begin{proposition}[Extension-exact conservative detection]\label{prop:detect}
Suppose genuine cartesian lifts exist over all set maps and all satisfy L4,
the content bodies $\eta^L,\mu^L$ are
invertible, and $d$ reflects isomorphisms of vertical arrows. Then L5 follows
from L1--L4 and cartesianness of $\eta^T,\mu^T$.
\end{proposition}
\begin{proof}
For a cartesian pair $(u:Q'\to Q,v:R'\to R)$ over $f$ and $du$, the two
display squares and preservation of pullbacks by $T$ identify
\[
 dR'\cong dR\times_{TTB}TTX.
\]
The structure map to $TTB$ is $T\nu_Q\,\nu_R$. Cartesianness of $\mu^T$
identifies this with $dR\times_{TB}TX$, now using
$\mu^T_B T\nu_Q\nu_R$. Invertible $\mu^L$ and~\eqref{eq:arity} identify
these with the extension display of the accumulated arrow. Factor that arrow
through the chosen cartesian lift of $Q\diamondc R$. Its vertical comparison
has bijective $d$-image because both sources solve the same pullback with the
same projections. Vertical conservativity makes it invertible. The unit
argument uses invertible $\eta^L$ and cartesian $\eta^T$ in the same way.
\end{proof}
The main theorem does not impose these additional detection hypotheses.
Without extension-exactness, the arity equations control only the images of
the bodies; the proof above would be invalid. Without vertical conservativity,
set bijectivity does not establish a native grade isomorphism.

\section{Construction of the full interaction}\label{sec:law}

\subsection{Generic crossing and extension body}
Let $P\in\K_B$ have extension projection $\pi:E=e_BP\to B$, and let
$Q\in\Lcat_B$. For $y\in dQ$, write
$\nu_Q(y)=(s,(b_i)_{i\in E_s})$. Define
\begin{equation}\label{eq:crossing}
 \widehat P_y=\boxtimes_{i\in E_s}P_{b_i},\qquad
 \widehat Q=\pi^*Q.
\end{equation}
The carrier of $\widehat P$ is canonically
\begin{equation}\label{eq:H}
 H_P(Q)=e_{dQ}\widehat P
 \cong dQ\times_{TB}TE.
\end{equation}
Its projection to $dQ$ is a pullback of $T\pi$, hence belongs to $\J$;
therefore $\widehat P$ is an admitted context grade. Display supplies the body
\[
 \chi_{P,Q}:H_P(Q)\xrightarrow{\sim}d(\pi^*Q),
 \qquad \chi_{P,Q}=(c^{\mathrm{disp}}_{\pi,Q})^{-1}.
\]
This body preserves both the original $y$ and the decorated $T$-label in $TE$.
Independent occurrences keep their addresses even when several labels $b_i$
are equal. Nullary occurrences give $J_\V$.

For a general representative
$(Q,f:A\to dQ)$ followed by $(P,g:eP\to C)$, reindex $\widehat P$ along
$f$, transport $\widehat Q$ along $g$, and use the body
\[
 e_A f^*\widehat P\longrightarrow H_P(Q)
 \xrightarrow{\chi_{P,Q}}d\widehat Q
 \xrightarrow{d\kappa_{g,\widehat Q}}d(g_!\widehat Q).
\]
This explicitly defines the output mixed representative, before descent
through the pseudocoend.

\subsection{Every native cell and pure endpoint map}
A native context arrow $P\to P'$ has carrier $h:E\to E'$ over $B$.
Its context component at $y$ is the occurrence tensor
of its actual aggregated coefficient maps, namely~\eqref{eq:tensormap}.
For its output grade, factor $c_{\pi,Q}:\pi^*Q\to Q$ through the cartesian
arrow $c_{\pi',Q}:\pi'^*Q\to Q$, giving an arrow over $h$. Then factor
through the opcartesian arrow over $h$ to obtain
\begin{equation}\label{eq:hmap}
 h_!\pi^*Q\longrightarrow\pi'^*Q.
\end{equation}
This comparison need not be invertible. Its functoriality follows by the two
universal properties; its coefficient functoriality follows from
\eqref{eq:tensormap}.

For a native content arrow $a:Q\to Q'$ over $B$, naturality of $\nu$ gives
$\nu_{Q'}da=\nu_Q$. Consequently the target occurrence set, shape and labels
are exactly the same at corresponding points; the context component is the
identity on the addressed coefficient objects. The grade component is
$\pi^*a$. This defines the action on all content arrows, not merely
isomorphisms. It does not assert that $\pi^*$ is faithful.

For both sorts of cells, the body equation is proved by composing the two
candidate maps with $d c_{\pi',Q'}$ and comparing their $TE'$-coordinates.
The first equality is the defining cartesian/opcartesian factorization; the
second is naturality of $\nu$. The target display pullback makes these two
coordinates jointly injective and gives the body equation. This use of
set-level uniqueness compares bodies only. Coefficient maps and native grade
maps have already been compared separately.

Arbitrary endpoint maps act by the given context reindexing and content
transport. Naturality with those maps follows from the same factorization
argument. In particular the native square formula~\eqref{eq:endpoints} is
preserved, not just its fixed-endpoint specialization.

\subsection{Pure-middle pseudobalance}
For a pure middle map $u:B\to B'$, the context reindexing square is
\[
 \begin{array}{ccc}
 E_u&\xrightarrow{v}&E'\\
 {\scriptstyle\pi_u}\downarrow&&\downarrow{\scriptstyle\pi'}\\
 B&\xrightarrow{u}&B'.
 \end{array}
\]
The output comparison between the two crossed representatives is precisely
\begin{equation}\label{eq:BC}
 \mathrm{BC}:v_!\pi_u^*Q\xrightarrow{\sim}\pi'^*u_!Q.
\end{equation}
The occurrence tensor contexts have the canonical addressed identification:
the pullback fibre over $b$ is the original fibre over $u(b)$. On coefficients
this is the identity, with the stipulated indexing comparisons.

To verify the body equation, a pair $(y,t)$ on the first side has target
coordinates
\[
 \bigl(d\kappa_{u,Q}(y),\ Tv(t)\bigr)
 \in d(u_!Q)\times_{TB'}TE'.
\]
The defining mate equation gives the same first coordinate on the second
side; naturality of $\nu$ gives the same second coordinate. Display at
$\pi'$ identifies the bodies. This is why no opcartesian preservation by
$d$ is needed. Naturality of BC follows by uniqueness for native grade maps;
context coefficient naturality is pointwise reindexing and
\eqref{eq:tensormap}. Identity and composition of balance generators are mate
pasting. Hence the construction respects every generator and relation of the
pseudocoend and yields a pseudonatural $\tau:N\odot M\Rightarrow M\odot N$.

\subsection{The four invertible structural modifications}
\paragraph{Context unit.}
For the constant singleton-$I$ context, $\pi\cong\id_B$,
$\pi^*Q\cong Q$, and $\bigotimes_{i\in E_s}I\cong I$.
The comparison is the canonical cleavage and tensor-unit isomorphism,
including $E_s=\varnothing$.

\paragraph{Content unit.}
The unit representative restricts the context along
$\eta^L_B:B\to dU_B$. At $y=\eta^L_B(b)$,
\eqref{eq:arity} gives the one-position polynomial unit shape, so
$\widehat P_y\cong P_b$. The output grade comparison is
$U_E\cong\pi^*U_B$ from~\eqref{eq:gradecomparisons}.
The lifted pair $(\eta^L_B(b),\eta^T_E(e))$ equals the image of
$\eta^L_E(e)$ by naturality and the display square. Thus the body equation
holds. No assertion that $dU_B\cong B$ is made: the comparison is on the
unit representative, not on all points of $dU_B$.

\paragraph{Context multiplication.}
Let $P$ over $B$ and $S$ over $E$ have carrier projections
$\pi:E\to B$, $\rho:F\to E$. The two grade outputs are
$(\pi\rho)^*Q$ and $\rho^*\pi^*Q$, canonically isomorphic by cleavage.
The context comparison is Proposition~\ref{prop:fubini}:
\[
 ((x_i,z_i))_i\longleftrightarrow((x_i)_i,(z_i)_i),\qquad
 \bigotimes_i(V_{x_i}\otimes W_{x_i,z_i})
 \cong(\bigotimes_iV_{x_i})\otimes(\bigotimes_iW_{x_i,z_i}).
\]
It is an augmented native isomorphism, natural even in fibre-merging maps.
Its body equation follows by display uniqueness with the same original
content value and the same decorated $T$-element.

\paragraph{Content multiplication.}
For $Q\in\Lcat_B$ and $R\in\Lcat_{dQ}$, the multiplied representative
restricts along $\mu^L_{Q,R}:dR\to d(Q\diamondc R)$.
A point $z\in dR$ specifies an outer shape and inner values $y_j\in dQ$.
Equation~\eqref{eq:arity} and cartesianness of $\mu^T$ identify the positions
of $\nu_{Q\diamondc R}(\mu^L(z))$ with addressed pairs $(j,i)$.
The context coefficient comparison is
\[
 \bigotimes_{\text{flattened leaves}}V_{\mathrm{leaf}}
 \cong\bigotimes_j\bigotimes_iV_{j,i},
\]
with that exact position bijection. The grade comparison is
\[
 \pi^*Q\diamondc j_\pi^*R\xrightarrow{\sim}\pi^*(Q\diamondc R),
 \qquad j_\pi=d c_{\pi,Q},
\]
from L5. Both body composites have first target display coordinate
\emph{$\mu^L_{Q,R}(z)$}, and have the same decorated flattened $T$-label.
The target display pullback therefore equates them. It need not remember
$z$ itself, and $\mu^L$ may be noninjective or nonsurjective.

All four components are invertible native comparisons by their stated
universal or structural properties. Their naturality follows from
Proposition~\ref{prop:fubini}, naturality of the polynomial transformations,
and cartesian uniqueness, not from an assumption that all native cells are
invertible.

\subsection{Full coherence without a faithful carrier shortcut}\label{sec:ten}
Use Gambino--Lobbia's Definition~4.1 and Appendix~A~\cite{GambinoLobbia} with
$S=N$, $T=M$, $(s,m)=(j_N,m_N)$ and $(t,n)=(j_M,m_M)$.
We verify C1--C10, including their two redundant diagrams. Their correspondence
with Marmolejo's named equations is
\[
\begin{array}{c|cccccccccc}
 \text{GL}&C1&C2&C3&C4&C5&C6&C7&C8&C9&C10\\\hline
 \text{coh}&4&2&5&3&1&6&9&7&\text{none}&8.
\end{array}
\]
An equality of mixed cells requires (a) equality of context carrier functions,
(b) equality of actual $\V$-coefficient maps, (c) equality of native $\Lcat$
grade maps, and (d) the comma body equation. None of (b) or (c) is inferred
from the set carrier.

For (a) and (b), work on the source representative of each diagram, retaining
all ancestor context choices and every addressed primitive weight occurrence.
When the diagram involves a content unit or multiplication, this representative
is restricted along the actual $\eta^L$ or $\mu^L$ body. We do not identify
that source with the entire extension of its accumulated grade. The structural
components have bijective context carriers, so each pushforward coefficient
fibre is singleton. Their coefficient maps are tensor associators, symmetries
and unitors on the addressed primitive objects. Both routes listed below
induce the same permutation of those occurrences, and symmetric monoidal
coherence identifies their actual coefficient maps.

\begin{description}[style=nextline]
\item[C1 content associativity on $NNNM$.]
Both routes substitute the same three-level polynomial tree. The monad
associativity equation identifies the flattened shape and the leaf map
$(k,j,i)\mapsto$ its original leaf. All context choices at those leaves remain
unchanged. The grade comparison is over the one-sided content associator.
\item[C2 content left unit.]
Restrict along the appropriate content-unit body. Both routes insert and
remove the same singleton-position layer, by $\mu^T(\eta^TT)=\id$.
The addressed coefficient word and every remaining context choice are fixed.
\item[C3 content multiplication against the context unit.]
Every primitive weight is $I$ and every choice is forced. Flattening the
content tree before or after removing these unit factors gives the same
canonical tensor-unit map, with the same multiplied content body.
\item[C4 both units.]
The polynomial unit has one position and the context unit has one choice.
Both pastings give the same $I$-unit map and the same pure body. The ambient
pure-action interchanger has this action as well.
\item[C5 both multiplications on $NNMM$.]
Address a primitive choice by its outer and inner content positions and its
outer and dependent inner context choices. One route flattens the two content
levels before splitting context pairs; the other splits inside each level
before flattening. Both send
$((j,i),(x_{ji},z_{ji}))$ to the same two addressed families. The coefficient
map is the same permutation from
$\bigotimes_{j,i}(V_{ji,x}\otimes W_{ji,x,z})$ to the separated outer and
inner tensor blocks. Proposition~\ref{prop:fubini} includes the necessary
interchange of actual native maps.
\item[C6 content unit against context multiplication.]
On the unit representative the one content position leaves the dependent pair
$(x,z)$ unchanged. One-position Fubini is this pair identification; both
coefficient maps are the same unit/reassociation map.
\item[C7 context associativity on $NMMM$.]
At each addressed leaf retain $(x,y,z)$, with the later fibres depending on
its preceding values. Both routes send $((x,y),z)$ and $(x,(y,z))$ to that
same triple and have the same tensor word $V_x\otimes W_{xy}\otimes Z_{xyz}$.
No later fibre is falsely treated as independent of its ancestors.
\item[C8 context right unit.]
The added singleton choice is forced. Removing it before or after crossing
preserves the same addressed history and deletes the same $I$ factor.
\item[C9 content right unit.]
On the representative restricted along the unit body, the trailing singleton
position is removed by $\mu^T(T\eta^T)=\id$. Both routes have the same
remaining occurrence word and body.
\item[C10 context left unit.]
An initial singleton context stage contributes only a forced choice and an
$I$ factor. Splitting then deleting it, or deleting before crossing, gives
the same addressed choices and coefficient map.
\end{description}

For naturality of these diagrams under arbitrary native coefficient cells,
expand a target fibre using~\eqref{eq:tensormap}. Both routes distribute,
apply the same addressed tensor network of aggregation maps, and regroup.
For dependent maps the continuation maps are applied before outer merging as
in~\eqref{eq:depmap}; subsequent maps are constant on the fibre being merged.
The detailed interchange in Lemma~\ref{lem:depfunctor} and the fixed-target
calculation in Proposition~\ref{prop:fubini} identify those networks.
Distributor naturality and rig coherence identify the surrounding structure.
Empty fibres use absorption naturality. This proves (b) for arbitrary native
cells, not only structural components.

For (c), compose the two native grade routes with the same target cartesian
arrow. One-sided content associativity/unit equations and the arity equations
identify the composites. Cartesian uniqueness identifies the vertical grade
maps. Where pure transport occurs, use the defining opcartesian factorization
and BC mate pasting first, then the same cartesian uniqueness. Thus every
parallel native map is retained. For (d), the body routes have the same
original target extension coordinate (after the appropriate $\eta^L$ or
$\mu^L$) and the same decorated polynomial coordinate. Display uniqueness
identifies them.

Finally, ambient tensor associators flatten only the list of balanced
representatives, and Gray interchangers compare natural maps in independent
slots. Their context coefficient actions are precisely the addressed
reassociations and exchanges just checked; their grade actions are the
factorization comparisons just checked. Lemma~\ref{lem:normalization} is fully
faithful on each hom category and therefore transports these equalities back
to the actual relative tensors. This proves all full law diagrams and their
endpoint naturality, not a strict-tensor surrogate.

\begin{corollary}\label{cor:law}
The sufficient conclusion of Theorem~\ref{thm:main} holds. The multiplication
\eqref{eq:composite} and its unit define the inherited composite relative
pseudopromonad on the original comma hom categories.
\end{corollary}

\subsection{Concrete mixed composition}
For $F=(P,Q,f:eP\to dQ):A\to B$ and
$G=(D,S,g:E=eD\to dS):B\to C$, cross $(Q,D)$ to obtain
$\widehat D\in\K_{dQ}$, $\widehat Q=\pi_D^*Q\in\Lcat_E$ and
$\chi:e\widehat D\to d\widehat Q$. The composite has context and content
\[
 P\starc f^*\widehat D,\qquad S\diamondc g_!\widehat Q,
\]
and body
\[
 e(P\starc f^*\widehat D)\cong e(f^*\widehat D)
 \longrightarrow e\widehat D\xrightarrow{\chi}d\widehat Q
 \xrightarrow{d\kappa_{g,\widehat Q}}d(g_!\widehat Q)
 \xrightarrow{\mu^L_{S,g_!\widehat Q}}d(S\diamondc g_!\widehat Q).
\]
The order $S\diamondc g_!\widehat Q$ matters. This formula retains the whole
output grade, not only points reached by the displayed body. It requires no
inverse of $\eta^L$ or $\mu^L$; the arbitrary map $g$ uses covariant transport.
The preceding law proof supplies its cell action, endpoint naturality and
associativity/unit coherence.

\subsection{Audit of every relative use}
Generic crossing uses display only at $\pi\in\J$. Context multiplication uses
$\pi,\rho,\pi\rho\in\J$. Content multiplication uses L5 over $\pi$ and
reindexes its second grade along $j_\pi\in\J$, a pullback of $T\pi$.
Context units use identities. Content units use L5 only over $\pi$ and restrict
along $\eta^L$. Pure-middle balance uses L6 on the context pullback square,
whose vertical arrows lie in $\J$. Arbitrary native carrier maps
$h:E\to E'$ factor through the genuine target $\pi'$-cartesian arrow and
then use the global opcartesian lift over $h$ in~\eqref{eq:hmap}. No $h^*$
is used. Body uniqueness uses target display at $\pi'$, not display at $h$. Native content cells and endpoints use $\nu$ naturality and
the supplied one-sided transport. Iteration invokes only the stated pullback
and composition closures. No all-map display or all-map cartesian-algebra
condition is hidden in this proof.

\section{Marked recognition and the exact necessary scope}\label{sec:recognition}

Let $\D$ be the class of carrier projections realized by admitted augmented
context families, up to isomorphism. Always $\D\subseteq\J$ and its maps have
finite fibres. Necessity cannot test a projection unavailable to the context
doctrine. For example normalized discrete nonzero semirings have no empty
fibres, while terminal or linear rigs admit them.

\begin{definition}[Marked exhaustive occurrence class]\label{def:marked}
Fix the one-sided data, global opfibration and genuine $\J$-cartesian lifts,
natural arity equations and the
context doctrine, without assuming L4--L6. A realization belongs to the marked
class when its generic crossed normal form has:
\begin{enumerate}[label=\textup{(\roman*)}]
\item output grade $\pi^*Q$ with its original cartesian arrow
$c_{\pi,Q}:\pi^*Q\to Q$ retained;
\item the complete context index set
$H_P(Q)=dQ\times_{TB}TE$, with all addressed choices, and an invertible
crossing body $\chi:H_P(Q)\to d(\pi^*Q)$ preserving both projections;
\item unit and multiplication grade comparisons preserving the original
arrows to $U_B$ and $Q\diamondc R$, together with the addressed context
coordinates;
\item pure-middle comparisons preserving the original commutative
cartesian/opcartesian lift square and its factorization equations.
\end{enumerate}
The last two clauses name independently available projections and lift arrows.
They do not supply a mixed comparison or assert that it is the canonical
cartesian comparison or BC mate.
\end{definition}

\begin{theorem}[Forced content exactness]\label{thm:recognition}
An interaction in Definition~\ref{def:marked} forces L4 and L5 along $\D$ and
L6 on the tested context reindexing squares. All its marked grade comparisons
are the canonical ones. Consequently the sufficient $\J$-criterion is an
if-and-only-if within this class when an independent richness argument gives
$\D=\J$.
\end{theorem}
\begin{proof}
For display, the canonical comparison $c^{\mathrm{disp}}_{\pi,Q}$ is defined
by naturality even before it is known invertible. Preservation of the two
marks gives $c^{\mathrm{disp}}_{\pi,Q}\chi=\id$. Invertibility of $\chi$
therefore gives $c^{\mathrm{disp}}_{\pi,Q}=\chi^{-1}$.

For the unit, factor the actual arrow $U_E\to U_B$ over $\pi$ through
$c_{\pi,U_B}$, obtaining $k_U:U_E\to\pi^*U_B$. The original-projection mark
makes the law's unit-grade component have the same composite with
$c_{\pi,U_B}$. Cartesian uniqueness forces it to equal $k_U$. Its
invertibility says exactly that the original unit arrow is cartesian.

For accumulation apply the independently defined accumulation to the cartesian
pair $(c_{\pi,Q},c_{j_\pi,R})$. Factor the result through the target cartesian
lift to obtain
\[
 k_a:\pi^*Q\diamondc j_\pi^*R\longrightarrow\pi^*(Q\diamondc R).
\]
The multiplication comparison preserves the original projection to
$Q\diamondc R$, so uniqueness forces its grade component to be $k_a$.
Its invertibility is equivalent to the accumulated arrow being cartesian.
This argument uses the generic source extension $dR$ and its actual body
$\mu^L$; it does not require $\mu^L$ invertible or $d$ faithful.

For pure balance, factor the marked comparison through the specified
opcartesian arrow over $v$ and the specified cartesian arrow over $\pi'$.
The projection mark says its composite is the original commuting lift square.
The two universal properties determine that factorization uniquely, so the
comparison is the canonical mate. Invertibility therefore forces BC for that
square. Thus neither canonicality nor the desired invertible comparison was
smuggled into the mark.

These comparisons and their equations are retained through the full comma
normalization, rather than an unmarked quotient. All tests are made at
realizable context projections. With $\D=\J$, they give the stated necessary
hypotheses, and Corollary~\ref{cor:law} supplies the converse.
\end{proof}

\begin{remark}[What a generic body alone proves]
Without exhaustive invertibility, preservation of the two coordinates says
only $c^{\mathrm{disp}}\chi=\id$: the comparison has a section. It need not
be injective. Invertible pseudolaw modifications do not make every crossing
body invertible. Accordingly this recognition theorem is not a theorem about
all unmarked distributive laws, correlated choices, quotient witnesses or
lax alternative normal forms.
\end{remark}

\subsection{Finite support removes one gap and no others}
\begin{proposition}[All finite-fibre display tests suffice]\label{prop:finitesupport}
In addition to the main theorem's relative lift existence, assume here that
$p:\Lcat\to\Set$ is a global fibration. Let $d:\Lcat\to\Set$ and
$\nu:d\Rightarrow Tp$ natural, with $T$ finitary polynomial. If every
cartesian lift over a finite-fibre map has a display pullback, then every
cartesian lift over every set map does. No continuity assumption on $d$ or
$\Lcat$ is needed.
\end{proposition}
\begin{proof}
Fix $f:X\to B$, $Q\in\Lcat_B$, and
$(y,t)\in dQ\times_{TB}TX$. Choose a finite subset $S\subseteq X$ supporting
$t$, with inclusion $i:S\hookrightarrow X$ and $t_S\in TS$ mapping to $t$.
Both $i$ and $fi$ have finite fibres. Display for $fi$ gives a unique
$w\in d((fi)^*Q)$ over $(y,t_S)$. Cartesian composition identifies this grade
with $i^*(f^*Q)$. Map $w$ to $z\in d(f^*Q)$; its coordinates are $(y,t)$.
This proves surjectivity.

If $z'$ is another lift, apply display for $i$ to the grade $f^*Q$. The pair
$(z',t_S)$ has a unique lift $w'$ in $d(i^*(f^*Q))$. In the identified grade
$(fi)^*Q$ its coordinates are $(y,t_S)$, hence $w'=w$ by $fi$-display
uniqueness. Thus $z'=z$. The canonical comparison is bijective.
\end{proof}
This proves only display exactness. It does not promote cartesian preservation
of accumulation or BC to all maps: those are native grade-isomorphism
statements, and $d$ need not reflect them. It also does not apply to tests only
on finite-fibre \emph{surjections}.

\begin{example}[A genuine contentad separating surjective tests]
Let $G(B)=\varnothing$ for $B=\varnothing$ and $G(B)=1$ otherwise.
The endofunctors $I=\id$ and $G$ form a strict two-object composition monoid,
$IG=GI=GG=G$. Give the constant grade category these objects and only identity
arrows, with extension $d_B(I)=B$, $d_B(G)=G(B)$, unit $I$, and accumulation
by functor composition. This is an actual one-sided contentad: its module
composition is functorial substitution and the monoid identities give its
pseudo-units and associativity. Constant reindexing is a bifibration, with
trivial BC and cartesian unit and accumulation.

Take $T=\List$, with $\nu_I(b)=[b]$ and $\nu_G(*)=[]$. Naturality and the
arity equations follow from singleton/empty-list substitution. For any
surjection $f:X\to B$, the $G$ display is a pullback: either both bases are
empty, or both $G$-values are singleton and the only list mapping to the empty
list is the empty list. The $I$ display is always cartesian. But for
$f:\varnothing\to1$, the $G$ comparison is
\[
 \varnothing=G(\varnothing)\longrightarrow
 G(1)\times_{\List(1)}\List(\varnothing)=1,
\]
not a bijection. Hence even valid contentad structure and every surjective
display test do not imply all-map display.
\end{example}

\subsection{Strict arity and native-cell boundaries}
For a native $a:Q\to Q'$ over $B$, the equality
$\nu_{Q'}da=\nu_Q$ preserves the entire polynomial shape and every label.
For $T=\List$, a binary-to-unary deletion cannot satisfy it, already at $B=1$.
For constant polynomial grades with cartesian display, the backward position
maps to the displayed shape are bijections; a native map commuting with the
display therefore has bijective position maps too. This theorem permits shape
merges with identical displayed arity, not arbitrary affine occurrence
deletion. Terminal augmentation makes context discard available, but does not
remove this separate strict-display restriction.

\section{Conservative controls and parameter changes}\label{sec:controls}

\subsection{Exact one-sided inclusions and trivializations}
The pure-content inclusion sends $(Q,f:A\to dQ)$ to
$(J_\V^A,Q,f)$, using $e_AJ_\V^A\cong A$. By
Lemma~\ref{lem:terminal} its context component in any native cell is uniquely
specified, so its native hom sets recover exactly those of $N$. The same
argument applies over arbitrary endpoints, since reindexing the singleton
unit gives the singleton unit. It is locally fully faithful.

The pure-context inclusion sends $(P,f:eP\to B)$ to
$(P,U_B,\eta^L_B f)$. Since $\eta^T_B$ is injective and
$\nu_{U_B}\eta^L_B=\eta^T_B$, the map $\eta^L_B$ is injective. Thus, if
$\End_{\Lcat_B}(U_B)=\{\id\}$, the mixed body equation cancels
$\eta^L_B$ and recovers exactly the original context body equation. This
proves local full faithfulness. For changing endpoints use the canonical
one-sided transport isomorphism $v_!U_B\cong U_{B'}$; grade maps to $U_{B'}$
then correspond to its endomorphisms, and naturality of $\eta^L$ gives the
same cancellation in~\eqref{eq:endpoints}.

Conversely, take input interface $A=\varnothing$. Its context grade category
is terminal and its extension is empty, so $M(\varnothing,B)$ has one object
and one endomorphism. In its mixed image every
$\beta\in\End_{\Lcat_B}(U_B)$ gives a distinct endomorphism, since the body
equation on the empty set is vacuous. Full faithfulness at all interfaces
therefore forces $\beta=\id$. This test remains available even when admitted
context projections are surjective: it uses $\id_\varnothing$, not an empty
fibre over a nonempty base.

Both inclusions are strong. Their multiplication comparisons are the
context-unit and content-unit crossing modifications, respectively; the latter
is restricted along $\eta^L$ exactly as above. Their units and associativity
coherence are the corresponding unit specializations of C1--C10. Thus strong
preservation is proved separately from local full faithfulness.

Identity-opposite recovery is different. Restrict contexts to the closed
subdoctrine consisting only of the singleton $J_\V$ at each input. It has
identity extension and gives the pure identity module, so mixing recovers the
original content construction. Conversely take the genuinely trivial content
$\Lcat_B=1$, $d_B(*)=B$, $T=\id$, with identity arity and bodies. Mixing
recovers the context construction. The terminal rig itself gives $\FinSet$
contexts, not identity contexts. Receipt content with $T=\id$ still has
nontrivial receipt tags and is not identity content.

\subsection{Strong morphisms of independently specified inputs}
\begin{definition}[Parameter map]\label{def:parameters}
Between two inputs (with a common admitted carrier class, or an inclusion of
source admitted projections in the target class), a parameter map consists of:
\begin{enumerate}[label=\textup{(\roman*)}]
\item a strong symmetric-rig functor $F:\V\to\V'$;
\item a cartesian morphism $\theta:T\Rightarrow T'$ of finitary cartesian
polynomial monads;
\item a strong contentad morphism $H:\Lcat\to\Lcat'$ over $\Set$ preserving
the available cartesian lifts and all opcartesian arrows up to their canonical
comparisons;
\item a coherent extension isomorphism $\delta:d'H\Rightarrow d$, compatible
with the contentad bodies and satisfying
$\nu'_{HQ}=\theta_{pQ}\nu_Q\delta_Q$.
\end{enumerate}
All strengths here are one-sided structure maps. No mixed-law map is supplied.
\end{definition}
Apply $F$ to every coefficient map and to the global augmentation, using its
strong sum and unit isomorphisms. Carriers are unchanged, so admissibility of
projections is preserved. Apply $H$ to content grades and use $\delta^{-1}$ on
bodies. This defines a functor on all comma cells because the body equation is
conjugated by the natural isomorphism $\delta$.

For the crossing comparison, cartesian preservation gives the unique
projection-preserving isomorphism
\begin{equation}\label{eq:Hpull}
 H(\pi^*Q)\xrightarrow{\sim}\pi^*(HQ).
\end{equation}
Cartesianness of $\theta$ identifies the position sets of a source shape and
its target shape, and strong tensor preservation by $F$ identifies
$F(\bigotimes_iV_i)$ with $\bigotimes_iFV_i$ along that bijection. Strong sum
preservation and augmentation naturality make this an augmented-family
isomorphism. Its body equality follows from the defining projections of
\eqref{eq:Hpull}, $\delta$-naturality and the arity equation.

For the four modification compatibilities, both coefficient routes expand to
$F$ of the same addressed tensor shuffle or unit map and then apply its strong
structure maps. The strong symmetric-rig functor axioms identify these actual
maps, including the naturality of finite distribution. The content routes are
maps between the same cartesian lifts; after composition with their original
grade projections they agree by the strong contentad unit/accumulation laws.
Cartesian uniqueness identifies them. BC compatibility follows likewise from
preservation of the defining cartesian/opcartesian factorization. Thus the
induced mixed functor is strong for~\eqref{eq:composite}.

There is also coherence under composition of parameter maps, not merely a
separate comparison for each map. For two composable maps $(F_1,H_1,\theta_1)$
and $(F_2,H_2,\theta_2)$, the position bijection for
$\theta_2\theta_1$ is the composite of the two bijections, since polynomial
transformations are determined by their position maps. The two context
comparison paths are therefore equal by strong functor coherence for
$F_2F_1$. The two grade paths
\[
 H_2H_1(\pi^*Q)\longrightarrow\pi^*(H_2H_1Q)
\]
have the same composite with the original cartesian projection: both are
$H_2H_1(c_{\pi,Q})$, with the prescribed preservation comparisons inserted.
Uniqueness makes them equal. Their body paths agree by the composite
$\delta$-equation. Identity parameter maps give identity comparisons by the
same arguments. Triple-composition coherence follows from these equality
statements, so no new higher parameter coherence is postulated.

If $F$ and the fibre functors of $H$ are faithful, the induced mixed functor
is locally faithful: carriers are unchanged, and equality of the two native
components can be reflected independently. If both are fully faithful, it is
locally fully faithful: fullness lifts coefficient and content arrows;
faithfulness of $F$ reflects the augmentation equation, and $\delta$ and
naturality reflect the body equation after lifting. The same reasoning with
transport comparisons handles endpoint squares.

Mere fullness of $F$ does not suffice. The full strong rig functor
$\FinVect_k\to1$ has no augmented-family lift of the unique image carrier map
from $(\{*\},k,\id)$ to $(\{*\},0,0)$: its only possible linear resource map
is zero, whose target augmentation composite is not $\id_k$. Thus fullness
without reflection of the augmentation equation cannot establish local
fullness. Likewise no general faithfulness of the crossing itself is claimed;
reindexing $\pi^*$ may identify native content arrows.

\subsection{A content change that genuinely changes the interaction}
For receipt content, a cartesian monad morphism $\theta:T\to T'$ defines
\[
 H_B(q:Q\to TB)=(\theta_Bq:Q\to T'B).
\]
Tags and extension $B+Q$ are unchanged. Unit preservation is
$\theta\eta=\eta'$; accumulation preservation is the monad-morphism equation
$\theta\mu=\mu' T'\theta\,\theta_T$, together with naturality. Cartesian
preservation is the pullback isomorphism
\[
 Q\times_{TB}TE\cong Q\times_{T'B}T'E.
\]
Opcartesian preservation is naturality under postcomposition by base maps.
Thus this is a parameter map with $\delta=\id$.

Take $T$ to be the free planar binary-tree monad (variables at leaves), $T'$ the
nonempty-list monad, and $\theta$ the left-to-right leaf list. Substitution of
trees becomes list flattening, proving the monad equations. Every tree has
exactly its list's ordered leaf positions, and relabelling those leaves is
unique, so $\theta$ is cartesian. Distinct bracketings of three leaves acquire
the same list label. The induced receipt functor is faithful on existing tag
maps, but not full: a map between singleton tags with different tree labels
and equal leaf lists is forbidden before flattening and allowed afterwards.
This changes actual content labels and which native merges exist, not merely
coefficients or an unused annotation. Occurrence multiplicities remain
preserved; arbitrary dependence-erasing maps are not asserted to satisfy the
parameter hypotheses.

\section{Different inputs satisfying the same criterion}\label{sec:instances}

\subsection{Four resource doctrines}
\paragraph{Discrete semiring resources.}
For a commutative semiring $R$, regard its elements as a discrete category,
with $\oplus=+$ and $\otimes=\cdot$. An augmentation exists precisely when
$\sum_xr_x=1$. A native map over $f:X\to Y$ exists precisely when
\[
 \sum_{fx=y}r_x=s_y
\]
for every $y$, and is then unique. Thus the construction recovers exactly raw
normalized finite weighted witnesses, with all zero-weight points retained.
Dependent accumulation multiplies $r_xs_{xy}$ and occurrence tensor takes
products of weights. For $0\ne1$ no empty normalized fibre exists. With raw
zero points and chosen nonempty finite fibres, assign weight $1$ to one point
and $0$ to the rest; hence every finite-fibre surjection is realizable in the
usual classical set setting. The normalized $R=\mathbb Q_{\ge0}$ or
$\mathbb R_{\ge0}$ doctrine and $T=\List$ receipt content recover the
inherited probability/list result, not a second newly counted theorem.

\paragraph{Terminal rig.}
For $\V=1$, every coefficient and arrow is unique. The augmented-family
category is $\FinSet$ with all functions, including empty families. Dependent
accumulation is finite dependent sum and occurrence tensor is finite product.
Every finite-fibre projection is realizable. This is not the nonempty,
surjective-map specialization coming from normalized Boolean weights.

\paragraph{Linear resources and invisible native maps.}
Take $\V=\FinVect_k$, with direct sum and tensor over a nonzero field.
An augmented family has one linear functional $\bigoplus_xV_x\to k$.
Native arrows are arbitrary fibre-aggregating linear maps satisfying the
augmentation equation. Over $k=\mathbb Q$, the singleton carrier with
$V=k^2$ and $\epsilon(a,b)=a+b$ has distinct parallel endomorphisms
\[
 \id(a,b)=(a,b),\qquad p(a,b)=(a+b,0),\qquad s(a,b)=(b,a).
\]
All preserve $\epsilon$; $p$ is noninvertible. They have the same carrier.
Under occurrence tensor and accumulation they act as actual maps such as
$p\otimes\id$, so their retention is mathematically substantive.
The unit $(\{*\},k,\id)$ nevertheless has only its identity endomorphism.

For an outcome-dependent example let an outer family have carrier $\{0,1\}$,
resources $k,k^2$ and augmentation adding all three coordinates. After outcome
$0$ choose one continuation with resource $k$; after outcome $1$ choose two
continuations with resource $k,k$, each with total-sum augmentation. The
accumulated carrier has three points and resources $k,k^2,k^2$; its
augmentation adds all five coordinates. The outer map $\id,p$ accumulates to
$\id,p,p$. It is undetectable on carriers and is not a scalar-weight map.
Zero vector spaces and the unique $0\to k$ augmentation realize every finite
fibre, including empty fibres.

This tensor is not a categorical dependent product in disguise. For the
singleton augmented object $(k^2,\epsilon)$ its tensor square has dimension
$4$, whereas its categorical product in the singleton slice over $k$ is the
pullback $k^2\times_k k^2$, of dimension $3$. Hence no identification with a
categorical-product construction follows merely from the word ``tensor.''

\paragraph{Capability-vector resources.}
Fix a finite coordinate set $K$. The thin category on
$R_K=(\mathbb N\cup\{\infty\})^K$ has an arrow $r\to s$ exactly when
$r\ge s$ coordinatewise. Put $\oplus=\min$ coordinatewise, $0_\oplus=\infty$,
$\otimes=+$ with $\infty$ absorbing, and $I=0$. These operations are monotone,
commutative and associative; addition distributes over minimum, including
$\infty$, and all coherence holds by thinness. Every finite family has a
unique augmentation $\min_xr_x\to0$, including empty families.
A native map over $f$ exists exactly when
\[
 \min_{fx=y}r_x\ge s_y,
 \quad\text{equivalently}\quad r_x\ge s_{fx}\text{ for every }x.
\]
Empty fibres impose no constraint. Dependent accumulation and occurrence
tensor add resource vectors. This is a thin \emph{inequality} rig category,
not the discrete normalized tropical semiring: no equation $\min_xr_x=0$ is
imposed. It therefore gives a genuinely nonprobabilistic resource doctrine
within the same construction. No free copying is supplied: for a positive
finite $r$, an arrow from singleton resource $r$ to singleton resource $r+r$
would require the false inequality $r\ge2r$.

\subsection{Indexed receipts with complete independent verification}
For a finitary cartesian polynomial monad $T$, let
$\Lcat_B=\Set/(TB)$. A grade is $q:Q\to TB$; a native map is a tag function
$a:Q\to Q'$ with $q'a=q$. Covariant transport along $f:B\to B'$ postcomposes
$q$ with $Tf$ and leaves tags unchanged. Cartesian transport is pullback
$Q\times_{TB}TX$ along $Tf$. These give a bifibration, and pullback
Beck--Chevalley follows from its universal property.
Set
\[
 dQ=B+Q,\quad \nu_Q=[\eta^T_B,q],\quad U_B=(\varnothing\to TB),
 \quad \eta^L_B:B\xrightarrow{\sim}B+\varnothing.
\]
For $r:R\to T(B+Q)$, accumulation has tag set $Q+R$ and label
\begin{equation}\label{eq:receiptmult}
 Q+R\xrightarrow{[q,\,\mu^T_B T\nu_Q r]}TB.
\end{equation}
The body $(B+Q)+R\to B+(Q+R)$ is reassociation. On composable-pair arrows it
is the coproduct of the actual tag maps; naturality of $\eta^T,\mu^T$ proves
its label equation. Covariant transport preserves the empty unit. Substitution
of three tag stages preserves the original stage tags, and its labels agree
by the monad associativity equation; the unit equations follow from the monad
unit equations. Coproduct associativity and unit coherence prove all one-sided
contentad laws on native arrows and endpoints. Thus this is not the invalid
slice-domain recipe of Remark~\ref{rem:slicefalse}.

For display, extensivity and cartesianness of $\eta^T$ give
\[
 (B+Q)\times_{TB}TX
 \cong X+(Q\times_{TB}TX)=d(f^*Q).
\]
The unit's cartesian preservation follows from the same calculation. For
accumulation, the first $Q$ summand of the pulled-back label gives $f^*Q$.
For the $R$ summand use
\[
 d(f^*Q)\cong dQ\times_{TB}TX,
 \quad
 T(d(f^*Q))\cong T(dQ)\times_{TTB}TTX,
\]
then the cartesian square of $\mu^T$ at $f$ to identify
$R\times_{T(dQ)}T(d(f^*Q))$ with the pullback of the flattened $R$ label
along $Tf$. This is exactly the second summand of $f^*(Q\diamondc R)$.
The isomorphism preserves both original tag projections, proving L5 on the
actual grade maps. All these checks hold for all set maps and hence for any
admissible $\J$. The unit grade is empty and has a unique endomorphism.
This verifies every independent hypothesis and both one-sided full-faithfulness
controls of the main theorem.

\subsection{Constant ordered polynomial content not generated by receipts}
Let $\mathcal C=\Set^{\mathbb N}$ be the category of ordered polynomial
collections $C=(C_n)_{n\ge0}$, with native maps the component functions.
Write
\[
 F_C(B)=\coprod_{n\ge0}C_n\times B^n.
\]
Let $\Lcat=\Set\times\mathcal C$ with constant reindexing and transport, and
$d(B,C)=F_C(B)$. A total arrow $(f,a):(B,C)\to(B',C')$ applies $a$ to its
shape and $f$ to every ordered position. Define $\nu$ to forget the shape and
retain the ordered list of labels, taking $T=\List$.

The unit collection has one unary shape and no others. For $C\in\Lcat_B$
and $D\in\Lcat_{F_CB}$ define $C\diamondc D$ to represent
$F_D(F_C B)$. A shape is an outer $D_m$ shape together with one $C_{n_j}$ shape
at each ordered outer position; its positions are the ordered concatenation
of the $n_j$ inner positions. Native maps act on all these shape components.
The body $F_D(F_CB)\to F_{C\diamondc D}B$ is the canonical substitution
bijection. Nested planar trees identify both triple substitutions with the
same fully addressed tree; these bijections are natural in every shape map.
Their pentagon and units are literal tree reassociation and deletion of unary
unit nodes. Hence this is a valid constant contentad with all native shape
merges, not only its core. Its module composition is ordinary functor
substitution in the execution order specified by the one-sided contentad.

Display at $f:X\to B$ is a pullback because a shape $c\in C_n$ is fixed and
a length-$n$ list over $X$ gives the unique relabelled point over that shape.
Cartesian arrows in the constant fibration have isomorphism grade components;
unit and substitution preserve those isomorphisms. BC is identity on grades.
The arity equations are singleton and ordered flattening. This verifies the
same L1--L6 as receipts, but its extension is polynomial evaluation rather
than $B+Q$, and grades are not endpoint-labelled receipt sets.

For example, two distinct binary shapes may be merged to one binary shape by
a noninvertible native arrow. This changes actual extension tags and is
retained by the theorem. A binary shape cannot be changed into a unary shape
in this category; arbitrary affine deletion or arbitrary reordering relative
to the fixed ordered-list display is not included. The unit collection has a
unique endomorphism, so both embeddings satisfy the exact control.

\subsection{A test with noninvertible content bodies}
For any cartesian polynomial monad $T$, take terminal grade fibres,
$d_B(*)=TB$, and the total extension action $Tf$. The contentad is the ordinary
monad regarded through its comma module: $N(A,B)=\Set(A,TB)$ with monad
substitution. Its unit and accumulation bodies are $\eta^T$ and $\mu^T$,
which need not be isomorphisms. Set $\nu=\id_T$. Its display square is
trivially a pullback, cartesian preservation is automatic in terminal fibres,
and BC is trivial. This legitimate instance demonstrates why the proof kept
$\eta^L$ and $\mu^L$ in every representative restriction.

\section{Contrasting operational support without application overclaims}\label{sec:fixtures}

The following small terminating syntax tests genuinely result-dependent
continuations in the same structural theorem. It is not the full O4 language,
an O2 external interface tensor, an O5 interference theorem, or a completed O6
application. In particular $\otimes$ above is an internal resource operation,
and $\boxtimes$ is independent occurrence combination. Neither proves a
physical scheduler-aware parallel composition.

\subsection{Finite probabilistic protocols}
Give a primitive at each input a finite declared outcome set and normalized
nonnegative rational weights, retaining zero-weight outcomes. The syntax has
\emph{skip}, dependent sequencing $P\mathbin{;}Q_x$ after the actual outcome
$x$, and a finite addressed batch of independent occurrences. Interpret skip
by the singleton of weight $1$, sequencing by carrier pairs $(x,y)$ with
weight $p(x)q_x(y)$, and batch by tuples with product weight.
Induction on syntax proves that the interpretation has exactly the declared
completed histories and their product weights. In the dependent case
$\sum_{x,y}p(x)q_x(y)=\sum_xp(x)=1$; batch normalization is finite
distributivity. Thus no history quotient or almost-sure identification enters
the interpretation. The general Fubini comparison is precisely the
history-preserving bijection between batching sequenced jobs and sequencing
an independent batch with its outcome-dependent continuation batch.

For a small exact fixture, job $A$ has outer outcomes $x,z$ of weights
$1/2,1/2$. After $x$, outcomes $u,v$ have weights $1/3,2/3$; after $z$ only
$w$ occurs with weight $1$. Completed histories have weights
\[
 (x,u):1/6,\qquad(x,v):1/3,\qquad(z,w):1/2.
\]
Let $B$ be a deterministic two-stage job. A receipt labelled $[A,A,B]$ has
nine distinct addressed histories with product weights, summing to $1$.
The two $A$ positions are independent declared occurrences, not one sampled
outcome copied into two coordinates. A genuinely shared-sample model would
need a different declared dependence structure. All equations here are exact
rational equalities and are checked in the companion regression script.

This is the inherited normalized protocol mechanism instantiated within the
new criterion. It supplies structural support for probabilistic applications,
not an observational or physical adequacy theorem for an adaptive laboratory.

\subsection{Finite consumable-cost protocols}
Use the capability-vector rig with coordinates $K$. A primitive protocol at
state $a$ has finite declared outcomes $X_a$ and exact costs
$c(a,x)\in(\mathbb N\cup\{\infty\})^K$. An empty outcome set means no
completed execution, not successful zero-cost execution. The syntax is
\begin{itemize}
\item a primitive outcome table;
\item skip, with one completed trace of cost $0$;
\item $\operatorname{then}(P,Q)$, with continuation $Q_x$ selected after the
actual trace $x$ of $P$;
\item $\operatorname{batch}(P_i:i\in E)$, with a finite set of distinct
addressed occurrences and the full product of their trace sets.
\end{itemize}
Interpretation sends primitives to their augmented cost families, skip to
$J_\V$, then to $\starc$, and batch to $\boxtimes$.

\begin{proposition}[Exact trace and cost interpretation]\label{prop:cost}
For every expression the interpreted carrier is exactly its finite set of
completed declared traces, and its resource at each trace is exactly the
operationally specified cumulative cost. The dependent Fubini isomorphism
preserves each addressed trace and its cost.
\end{proposition}
\begin{proof}
Induct on syntax. Primitive and skip cases are definitions. For then,
accumulation has carrier $\coprod_x|Q_x|$ and resource
$c_P(x)+c_{Q_x}(y)$. For batch its carrier is the product and its resource is
$\sum_i c_{P_i}(x_i)$. The global augmentation exists since every vector is
at least $0$. Fubini sends $((x_i,y_i))$ to $((x_i),(y_i))$ and both costs
are $\sum_i(c_{P_i}(x_i)+c_{Q_{i,x_i}}(y_i))$, by finite associativity and
commutativity. This is exactly Proposition~\ref{prop:fubini} in this rig.
\end{proof}
Batch assumes independently admissible choices, records cumulative consumption
rather than peak occupancy, and does not encode capacity contention, hidden
joint constraints, concurrency duration or physical interleavings. The
isomorphism is a denotational trace/cost equivalence, not permission to reorder
physical side effects.

Use actual receipt content with $T=\List$ and job labels $\{A,B\}$. Give
three receipt tags the lists $[A]$, $[A,B]$, $[A,A,B]$. The crossing retains
each original tag and its addressed tuple, sums the costs over those positions,
and its multiplication comparisons are the trace/cost Fubini and leaf
flattening just proved. This is therefore a content instance of the main law,
not only an interpretation of its resource objects.

For coordinates (fuel, material), let $A$ have outer outcome $x$ of cost
$(1,2)$ with continuations $u:(2,0)$ and $v:(1,1)$, and outer outcome $z$ of
cost $(2,1)$ with continuation $w:(0,2)$. Let $B$ have one outer outcome of
cost $(1,0)$ and one continuation of cost $(0,1)$. Then
\[
 A:\ (x,u):(3,2),\quad(x,v):(2,3),\quad(z,w):(2,3),\qquad B:(1,1).
\]
The receipt $[A]$ has these three costs. The receipt $[A,B]$ has costs
$(4,3),(3,4),(3,4)$. The receipt $[A,A,B]$ has nine histories, with costs
$(7,5)$ once, $(6,6)$ four times and $(5,7)$ four times. These multiplicities
are not probabilities and are not divided by nine. Crossing before or after
accumulation yields identical addressed histories and costs.

\begin{proposition}[Native per-outcome budget certificate]\label{prop:budget}
For a cost family $D=(X,d_x)$ and a budget vector $b$, let $B_b(X)$ have the
same carrier and constant resource $b$. An identity-carrier native arrow
$B_b(X)\to D$ exists if and only if $b\ge d_x$ for every $x\in X$.
For nonempty finite $X$, the least such $b$ is the coordinatewise maximum of
the $d_x$.
\end{proposition}
\begin{proof}
Each identity-carrier fibre is singleton, so the native arrow condition is
exactly the stated inequality at that point. Taking coordinatewise maxima
is equivalent to all these inequalities. For $X=\varnothing$ the condition
is vacuous and proves nothing about existence of a completed run.
\end{proof}
The rig's additive minimum is not a worst-case upper bound. The three receipts
above admit budgets $(3,3)$, $(4,4)$ and $(7,7)$; $(4,3)$ fails for $[A,B]$
because two histories cost $(3,4)$. Its minimum augmentation uses $(3,3)$,
which is not a universally sufficient budget.

Certificates compose by actual native maps. Given $b\ge c_P(x)$ and
$b_x\ge c_{Q_x}(y)$, dependent accumulation gives source resource $b+b_x$
at $(x,y)$. A uniform budget requires separately $B\ge b+b_x$ for every $x$;
composing that identity-carrier arrow gives the uniform certificate.
Independent tensors of certificates add their budgets. Positive finite
resources admit no free-copy arrow $r\to2r$. No undo or recovery operation
is postulated.

This proves nonprobabilistic, no-recovery structural support with operational
induction. It does not establish uncertain/discovered cost bounds, physical
safety, adversarial scheduling, freshness or schema transitions, typed
admissible deployment operations, or a complete application interpretation.
Those stronger motivating applications remain later work under the expanded
contract, rather than being silently replaced by this fixture.

\section{A distinct display and codisplay branch}\label{sec:pbpo}

This section records a separate independently derived positive control and
bounded obstruction for the earlier pullback/image/pushout recipe. It is not
joined by disjoint union to the occurrence theorem and advertised as one
universal criterion. The nonthin rig proof above does not follow from the
thinness argument below, and conversely these exactness results do not require
a polynomial occurrence display.

\subsection{Abelian bases and independently chosen Serre classes}
Let $\A$ be a universe-bounded abelian category. Choose independently full
replete Serre classes $\Uclass,\Wclass$: they contain $0$ and are closed under
subobjects, quotients and extensions. Displays at $A$ are actual monos
$D\hookrightarrow A$ with $A/D\in\Uclass$, with all maps over $A$.
Codisplays at $B$ are actual epis $B\twoheadrightarrow B/N$ with
$N\in\Wclass$, with all quotient maps under $B$. Retaining actual presentations
retains their canonical isomorphisms; passing to a subobject skeleton is an
equivalence, not a quotient by witness connectivity.

Pullback displays remain admitted because their cokernel embeds in the original
cokernel. Composite display cokernels are extensions of the two original
cokernels. Pushout codisplays have kernel $\im(vN)$, a quotient of $N$;
composite epi kernels are extensions. Therefore the one-sided context and
content structures exist independently, with pullback and pushout transport.
Their native categories are thin up to presentation isomorphism by mono and
epi cancellation, but contain noninvertible inclusions and coarsenings.

For a codisplay $q:B\twoheadrightarrow B/N$ followed by a display
$E\hookrightarrow B$, set
\[
 W=(E+N)/N\cong E/(E\cap N).
\]
The crossing is the display $W\hookrightarrow B/N$, the codisplay
$E\twoheadrightarrow E/(E\cap N)$, and the indicated canonical body
isomorphism. Its display cokernel is $B/(E+N)$, a quotient of $B/E$; its
codisplay kernel is $E\cap N$, a subobject of $N$. Both are admitted.
The comparison is derived by abelian exactness, not assumed as mixed data.

\subsection{A faithful relation presentation with all endpoint squares}
A mixed form is $(D\hookrightarrow A,B\twoheadrightarrow B/N,f:D\to B/N)$.
Associate the relation subobject
\begin{equation}\label{eq:relation}
 R=D\times_{B/N}B\hookrightarrow A\oplus B.
\end{equation}
Conversely, for $R\hookrightarrow A\oplus B$ let $D=\im(R\to A)$ and let
$N=\ker(R\to A)$, regarded as a subobject of $B$. The epi $R\to D$ is the
cokernel of $N$. The composite $R\to B\to B/N$ kills $N$ and therefore
descends uniquely to $f:D\to B/N$. The induced arrow
$R\to D\times_{B/N}B$ is an isomorphism: it is a morphism of the short exact
sequences with kernel $N$ and quotient $D$, identical at both ends. This proves
exact reconstruction, internally in any abelian category.

A fixed-endpoint native cell exists precisely when $D\subseteq D'$,
$N\subseteq N'$, and the body triangle commutes. This is equivalent to
$R\subseteq R'$: inclusion implies the two inclusions by image/kernel and the
body equation by epi descent; conversely the body equation defines the map of
pullbacks. Both native maps are unique by cancellation. Over arbitrary pure
endpoints $(u:A\to A',v:B\to B')$, the corresponding condition is
\[
 \im((u\oplus v)|_R)\subseteq R',
\]
equivalently $uD\subseteq D'$, $vN\subseteq N'$ and the quotient-body equation.
Hence~\eqref{eq:relation} is a fully faithful presentation on the actual
normal-form hom categories and endpoint squares. Its image consists exactly
of relations with $A/D\in\Uclass$ and $N\in\Wclass$.

\subsection{The original composition and its balance generator}
Compose the preceding form with
$(E\hookrightarrow B,C\twoheadrightarrow C/L,g:E\to C/L)$.
The original recipe forms $W=q(E)$, restricts to $D'=f^{-1}W$, and pushes out
\[
 P=W+_E(C/L).
\]
Since $E\twoheadrightarrow W$ is epi, $C/L\to P$ is epi and
\[
 P\cong(C/L)/\im(g|_{E\cap N}).
\]
The kernel of $C\to P$ is an extension of $L$ by that image, hence lies in
$\Wclass$. Also $D/D'$ embeds in $(B/N)/W=B/(E+N)$, a quotient of $B/E$;
$A/D'$ is an extension of $A/D$ by $D/D'$, hence belongs to $\Uclass$.

The abelian pushout has the cokernel presentation
\[
 P=(W\oplus C/L)/\im\bigl(E\xrightarrow{(q,-g)}W\oplus C/L\bigr).
\]
Its pullback comparison $E\to W\times_P(C/L)$ is epi: the pullback is the
kernel of the map $(w,z)\mapsto[w]-[z]$ to the displayed cokernel, and this
kernel is exactly the image of $(q,g)$. Thus the relation of the new form is
ordinary image-of-pullback relational composition. In element language an
equality in the pushout is witnessed by a single $e\in E$; additivity combines
all gluing witnesses into that one image. No difunctional closure is inserted.

For clarity, pure-middle balance is checked explicitly, not inferred only
from associativity. Let $R\hookrightarrow A\oplus B$ and
$S\hookrightarrow B'\oplus C$ be the two relations and $u:B\to B'$ pure.
Covariant output transport of $R$ is the regular image
$(\id_A\oplus u)_*R$, as the pushout quotient is $B'/\im(uN)$.
Contravariant input transport of $S$ is the inverse image
$(u\oplus\id_C)^*S$. The balanced representatives have equal composites:
\begin{equation}\label{eq:relbalance}
 S\circ(\id_A\oplus u)_*R
 =((u\oplus\id_C)^*S)\circ R.
\end{equation}
Indeed each side is the image in $A\oplus C$ of the same pullback condition
on $R$ and $S$, namely the equality of the $S$ input coordinate with $u$
applied to the $R$ output coordinate. Factoring through the intermediate
regular image does not change this final image, because regular epis are
pullback stable. This proves the balance generator internally; its native
naturality and identity/composition relations follow from functoriality of
images and pullbacks, then from the faithful presentation above.

Relational composition is associative: iterated image-of-pullback composition
has the image of the common triple-relation pullback, by the same stability of
regular epis. The diagonal relation is its identity. Crossing and the four
unit/multiplication normalizations all represent the same corresponding
relation by the preceding calculation and~\eqref{eq:relbalance}. Full
faithfulness gives their invertible native comparison. Every higher coherence
diagram commutes because between the fixed source and target forms with fixed
endpoints there is at most one native map. This uses genuine thinness proved
for this family; it erases no native noninvertible map and gives no shortcut
for the nonfaithful linear-resource case of the main theorem.

The one-sided embeddings are $N=0$ and $D=A$. The same presentation shows
local full faithfulness, and the composition calculation gives strong
preservation of their original pullback and pushout composition. An exact
additive functor carrying $\Uclass$ into $\Uclass'$ and $\Wclass$ into
$\Wclass'$ preserves finite limits, colimits and images, hence strongly
preserves the crossing and all its canonical comparisons. Faithfulness or
fullness of this functor is not automatic.

This covers module categories, abelian functor categories and chain-complex
categories. The two Serre classes may vary independently; for example torsion
supported at two independently chosen sets of primes in $\Ab$. The subsequent
display can depend on the actual preceding output by inverse image, so this
is a genuinely dependent indexed positive branch, though its native grades
are thin.

\subsection{A sharp bounded group recognition theorem}
\begin{theorem}\label{thm:groups}
Let $\mathcal G$ be a nonempty full subcategory of $\Grp$, closed under
subgroups, quotients and finite products. Use all subgroup displays and quotient
codisplays, arbitrary pure group homomorphisms, and the original
image/pullback/pushout crossing. The following are equivalent:
\begin{enumerate}[label=\textup{(\roman*)}]
\item its outputs respect pure-middle balance;
\item every group in $\mathcal G$ is abelian;
\item it extends to the retained-cell coherent law just proved, internally
with both Serre classes equal to $\mathcal G$.
\end{enumerate}
\end{theorem}
\begin{proof}
The implication (iii)$\Rightarrow$(i) is immediate. If (ii) holds, the full
subcategory is abelian: it contains zero, kernels and cokernels by the closure
assumptions, finite products are biproducts, and additive hom structures are
inherited by fullness. The standard abelian image/coimage isomorphism remains
in the subcategory. Taking both classes to be the whole internal abelian
category proves (iii). No closure under arbitrary ambient extensions in
$\Ab$ is required for this argument.

For (i)$\Rightarrow$(ii), fix $H,K\le G$. Use the quotient $H\to1$, the pure
inclusion $H\hookrightarrow G$, and the partial continuation with domain
$K\hookrightarrow G$ and identity body. Pushout first yields, after crossing,
\[
 K\longrightarrow K/(K\cap\langle\!\langle H\rangle\!\rangle_G),
\]
whereas pullback first yields
\[
 K\longrightarrow K/\langle\!\langle H\cap K\rangle\!\rangle_K.
\]
The source representatives are identified by pure balance, so their retained
quotient kernels must agree. For any
$x\in\langle\!\langle H\rangle\!\rangle_G$, take $K=\langle x\rangle$.
This $K$ is cyclic, so $H\cap K$ is already normal. The first kernel is all
of $K$, hence so is the second, giving $x\in H$. Therefore every subgroup of
every $G$ is normal.

Apply this conclusion in $G\times G$ to its diagonal. For every $g,h\in G$,
normality gives $(ghg^{-1},h)\in\Delta G$, and therefore $ghg^{-1}=h$.
Thus $G$ is abelian.
\end{proof}
For $G=S_3$, let $H=C_2$ be generated by a transposition and $K=C_3$ the
alternating subgroup. Their intersection is trivial while the normal closure
of $H$ is all of $S_3$. The two balance outputs are $K\to1$ and $K\to K$.
Hence exact Mal'cev or semi-abelian hypotheses do not suffice for this recipe.
The issue is normal closure generated outside the continuation domain, not
merely the failure of set relations to be difunctional.

The proof is constructive for any nonabelian $G$: in $G\times G$ choose
$H=\Delta G$ and a conjugate $x=(ghg^{-1},h)\notin H$; take $K=\langle x\rangle$.
The first quotient kills all of $K$, the second does not. This bounded theorem
does not exclude different crossings, restricted displays, larger normal
forms, lax laws or other nonabelian mechanisms.

\subsection{Disposition of the branch}
The local exactness fact used above is the regular-pushout condition
$E\twoheadrightarrow W\times_{W+_EZ}Z$ for one epi leg $E\twoheadrightarrow W$
and arbitrary $E\to Z$. It is stronger than merely having a pushout, and is
not the two-epi statement often available in exact Mal'cev categories.
The abelian proof establishes it by the displayed cokernel calculation.
We have not shown that this condition alone, in an arbitrary regular category,
reconstructs every native normal form or implies the complete law; the faithful
relation reconstruction is another essential proved step here.
The branch therefore has a bounded positive result, a sharp original-recipe
group obstruction, and an explicit remaining wider problem. It is neither an
unrestricted nonthin repair nor a claim that all nonabelian alternatives fail.

\section{Strongest source comparisons and mathematical increment}\label{sec:sources}

The assembly uses established polynomial, fibrational and distributivity
coherence mechanisms. It is not presented as the discovery of generic
coherence by uniqueness. The following comparisons identify actual overlap
and the remaining translation, rather than inferring novelty from different
terminology.

\paragraph{Higher-dimensional pullback laws and relative maps.}
Von Glehn~\cite{vonGlehn}, Proposition~5.1, already derives a full
pullback-driven pseudodistributive law. Theorem~7.1 treats fibred sums and
products in a locally cartesian closed base and Theorem~7.2 identifies the
polynomial double-category composite. Section~8 explicitly treats a chosen
pullback-stable display-map class and relative sums/products. Thus the relative
$\J$ organization and universal-property coherence have strong precedent.
Our arbitrary rig occurrence tensor need not be a categorical right-adjoint
dependent product, as the dimension-$4$ versus dimension-$3$ linear example
shows; our supplied $\Lcat$ need not be the free fibred-product doctrine.
A direct specialization would have to identify those structures and all native
maps, rather than just invoke that theorem's vocabulary.

\paragraph{Double clubs and nonflat cells.}
Garner's \emph{Double clubs}~\cite{DoubleClubs} is arXiv
\texttt{math/0606733}; \emph{Polycategories via pseudo-distributive laws}
~\cite{Polycategories} is \texttt{math/0606735}. They are distinct papers.
The former's cartesian-collection and evaluation-at-$1$ machinery supplies
coherence from horizontal pullback stability. The latter's Proposition~20
lifts through the relevant cartesian-collection slice; Proposition~38 and
Corollary~39 establish and use local monomorphism for its particular matching
construction to prove the ten axioms. These are genuine retained-cell results.
They do not say that a nonfaithful coefficient carrier such as
$\Aug(\FinVect_k)\to\FinSet$ is locally monic. The present coefficient proof
therefore checks resource maps directly. An equivalence with an appropriate
double-club presentation could turn much of the assembly into a corollary;
that equivalence is not presumed or ruled out.

\paragraph{Fibrational and polynomial mechanisms.}
Shulman~\cite{Shulman}, Theorem~14.2, constructs a framed bicategory from
frameable monoidal bifibration data, and Theorem~14.9 gives functoriality.
This is independent-hypothesis coherence machinery, not merely a Set-level
analogy. Our dependent content accumulation has not automatically been
identified with the external monoidal structure required there. Weber
~\cite{Weber} constructs polynomial composition by distributivity pullbacks
and its 2-categorical extension. The polynomial leaf and cartesian-mate
mechanisms used here are established; the task discharged here is their
explicit joint action on the supplied augmented coefficients and content
native cells. Laplaza's regular-expression qualification and the extra use of
native-arrow naturality were stated precisely in Section~\ref{sec:rig}.

\paragraph{The ordinary independent-input baseline.}
Fujii--Tsai--Montacute--Hasuo~\cite{FTMH}, Theorem~42, gives an independent
$W$-operadic/$W$-commutative criterion for an ordinary distributive law over
the underlying functor, and a full monad law with their stated additional
hypotheses. This includes the familiar commutative probabilistic shadow.
It is not stated as the indexed retained-grade, pure-pseudobalanced theorem
proved here. A surviving ordinary law after erasure does not by itself imply
an invertible retained structural comparison; this is a distinction of
specifications, not a contradiction of that baseline.

\paragraph{The abelian branch and factorization.}
Fong--Zanasi~\cite{FongZanasi}, Section~3.2 and Corollary~3.4, already give
abelian relation/corelation and span/cospan constructions; their assumptions
make an appropriate pushout/pullback comparison explicit, and their finite-set
nonexample warns against separate pullback/pushout existence alone.
The branch above makes the independently chosen Serre classes, native endpoint
presentation and bounded group obstruction explicit for this programme; no
priority is claimed for that packaging. Rosebrugh--Wood~\cite{RosebrughWood}
and \v{S}t\v{e}p\'an~\cite{Stepan} provide important factorization baselines.
Starting with a desired mixed faithful amalgam and then invoking a
factorization theorem would not answer existence; our relation reconstruction
and our occurrence construction are instead independently proved.

The defensible increment is a common explicit retained-cell assembly for the
stated independent $\V,\Lcat$ inputs, its bounded forced exactness, exact
controls and discriminating instances. The checked sources do not make it an
immediate corollary in this presentation, but that is not evidence of an
unfilled historical gap or a novelty certificate. No exhaustive priority
assessment is claimed.

\section{O1 crosswalk and exact remaining limits}\label{sec:acceptance}

The broader result now meets the mathematical O1 package at its explicit
relative augmented-rig/occurrence-display scope and is ready to present for
acceptance. Proof validity and the user's acceptance of this structural scope
remain distinct. The original O1 was a schematic family theorem, not a promise
of faithful interaction for arbitrary pairs. Conversely the receipt-only result
is not relabelled as the broader result: the independent contentad varies here.

\begin{longtable}{@{}p{.29\linewidth}p{.65\linewidth}@{}}
\toprule
O1 requirement&Evidence and exact status\\
\midrule\endhead
Independent entry test&Definition~\ref{def:content}, the rig axioms and the
ledger mention no mixed law. Receipts and constant polynomial content satisfy
the same tests through different independent presentations.\\
Full interaction and relative algebra&Sections~\ref{sec:rig} and
\ref{sec:law} construct all actual coefficient and grade maps, body equations,
endpoint actions and pure balance, four modifications and C1--C10;
\eqref{eq:composite} is the inherited composite theorem.\\
Substantive parameter change&Definition~\ref{def:parameters} and its proof
give strong functoriality and parameter-composition coherence. Tree-to-list
change alters actual content labels and admissible merges.\\
Conservative one-sided controls&Section~\ref{sec:controls} proves strong pure
content local full faithfulness, pure context local full faithfulness exactly
under unit rigidity, and the correct identity-opposite trivializations.\\
Bounded necessity beyond marked presentation recovery&Theorem
~\ref{thm:recognition} forces cartesian unit/accumulation and canonical BC
from original projection marks; exhaustive body exactness separately forces
display. It states the $\D$ versus $\J$ gap and
Proposition~\ref{prop:finitesupport} strengthens display alone.\\
Result-dependent continuations&The native dependent accumulation permits the
next grade to depend on the actual previous witness; both operational
inductions exercise it.\\
Contrasting O6 support&Section~\ref{sec:fixtures} proves normalized history
weights and nonprobabilistic trace/cost semantics inside the same criterion,
including exact native budget certificates without recovery. This is support,
not completed O6 acceptance.\\
Inherited probability/list control&Discrete semiring resources and list
receipts exactly recover the earlier raw weighted construction, including
zero points, native maps and labels. It is not counted twice.\\
Strongest-source comparison&Section~\ref{sec:sources} credits existing full
categorical and relative mechanisms, specifies non-reductions, and makes no
novelty or priority claim.\\
Original PB/PO disposition&Section~\ref{sec:pbpo} retains a distinct
abelian/Serre positive-control branch and exact original-recipe group iff;
arbitrary nonthin or nonabelian alternatives remain outside it.\\
\bottomrule
\end{longtable}

The following boundaries remain material, not footnotes that can be removed
by changing terminology.
\begin{enumerate}
\item The context family comes from coherent symmetric rigs and finite
witnesses. This is not a classification of arbitrary contextads. Infinite
occurrences require further existence or convergence data.
\item Content must already be a valid one-sided contentad. Its relative
exactness and strict polynomial arity conditions are substantive, and its
native maps need not include occurrence deletion or other arity changes.
\item Necessity is for exhaustive marked addressed normal forms and only
along realizable projections, unless the proved richness or finite-support
arguments apply. Arbitrary unmarked laws, correlations, quotient witnesses and
lax alternatives are not classified.
\item Retaining all native arrows does not imply that the crossing or a
parameter map is faithful. Faithfulness/fullness and unit rigidity have their
stated separate hypotheses. No set-level bijection is silently promoted to a
native grade isomorphism.
\item The occurrence tensor is not an O2 external interface tensor, arbitrary
copying operation or physical concurrency law. The fixtures do not prove
interference properties, physical resource safety, deployment typing,
freshness, schema transitions, incomplete-observation adequacy or recovery.
\item O2 structural integration, O3 comparisons including the QBS benchmark,
O4 language, O5 semantic soundness/refinement, the complete O6 application
sheets, and other expanded-contract obligations are not discharged by this
O1 result. Fuller motivating applications can proceed later as requested.
\end{enumerate}
Thus the whole expanded contract remains open. A future request for a wider
family can extend this explicitly stated result; it should not be satisfied
by silently broadening its quantifiers or silently shrinking its obligations.

\appendix
\section{Reproducible checks and their evidential limits}\label{app:checks}
The accompanying source-and-check bundle contains this TeX source, a README,
manifest and SHA-256 hashes, plus the exact standard-library regression scripts
and their results. It deliberately excludes primary-source PDFs and private
working notes. No TeX engine was run, by the authorized source-only workflow;
no PDF or visual-layout verification is claimed.

\begin{itemize}
\item \texttt{check\_rig\_native\_cells.py} checks exact rational linear
aggregation, dependent composition and occurrence/Fubini naturality, including
parallel maps invisible on carriers. It is a discriminator for the actual
coefficient proof, not a proof for all rigs.
\item \texttt{check\_capability\_fixture.py} verifies every addressed trace
and vector cost in the stated fixture, budget success and failure, absence of
free copying for the tested positive vectors, and all job words of lengths
zero through five.
\item \texttt{check\_probability\_fixture.py} uses exact rational arithmetic
to compare sequencing and occurrence batching and their total masses, including
the empty occurrence family and retained zero-weight histories.
\item \texttt{check\_group\_balance.py} enumerates actual finite permutations,
subgroups, normal closures and cosets for the $S_3$ example and diagonal group
witness. It verifies the unequal native quotient sizes, not the general iff.
\item \texttt{check\_tex\_source.py} checks balanced braces and environments,
unique labels, resolved internal references and bibliography keys, and expected
scope safeguards. Static success does not assert compilation or mathematical
correctness.
\end{itemize}
The general theorems are established by the proofs above. Tests are finite
regressions and transcription checks; they cannot replace the universal
coefficient, cartesian-uniqueness or recognition arguments.

\begin{thebibliography}{99}
\bibitem{Contract}
\emph{Expanded Research Contract Version 0.1}. Project TeX source,
1 November 2026, O1 and the fixed substantive minimum. Library ID
\texttt{}.

\bibitem{Receipt}
\emph{Indexed Receipt Interaction Theorem}. Project proof dossier,
1 November 2026. Scoped normalized-weighted/receipt predecessor, retained
unchanged; source filename \texttt{Indexed\_Receipt\_Interaction\_Theorem.tex}.

\bibitem{Unified}
\emph{One relative construction for context content and effects}. Project
TeX source, 1 November 2026. Library ID
\texttt{}.
Pseudocoend normalization, retained endpoint realization and supplied-law
composite machinery.

\bibitem{GambinoLobbia}
N. Gambino and G. Lobbia. \emph{On the formal theory of pseudomonads and
pseudodistributive laws}. Theory and Applications of Categories 37 (2021),
no. 2, 14--56. Definition 4.1 and Appendix A; arXiv:0907.1359v3.
\url{https://arxiv.org/abs/0907.1359v3}.

\bibitem{Marmolejo}
F. Marmolejo. \emph{Distributive laws for pseudomonads}. Theory and
Applications of Categories 5 (1999), no. 5, 91--147. Section 4 and Theorem 5.1.
\url{https://www.tac.mta.ca/tac/volumes/1999/n5/5-05abs.html}.

\bibitem{Laplaza}
M. L. Laplaza. \emph{Coherence for distributivity}. Lecture Notes in
Mathematics 281 (1972), 29--72. Regularity in Section 3 and Proposition 10.
\url{https://www.sas.rochester.edu/mth/sites/doug-ravenel/otherpapers/laplaza.pdf}.

\bibitem{vonGlehn}
T. von Glehn. \emph{Polynomials, fibrations and distributive laws}. Theory
and Applications of Categories 33 (2018), no. 36, 1111--1144.
Proposition 5.1, Theorems 7.1--7.2 and Section 8.
\url{https://tac.mta.ca/tac/volumes/33/36/33-36.pdf}.

\bibitem{DoubleClubs}
R. Garner. \emph{Double clubs}. Cahiers de Topologie et G\'eom\'etrie
Diff\'erentielle Cat\'egoriques 47 (2006), no. 4, 261--317.
arXiv:math/0606733. Definitions and propositions on cartesian collections and
horizontal pullback stability.
\url{https://www.numdam.org/item/CTGDC_2006__47_4_261_0.pdf}.

\bibitem{Polycategories}
R. Garner. \emph{Polycategories via pseudo-distributive laws}.
arXiv:math/0606735v3. Proposition 20, Proposition 38 and Corollary 39.
\url{https://arxiv.org/abs/math/0606735}.

\bibitem{Shulman}
M. Shulman. \emph{Framed bicategories and monoidal fibrations}.
Theorems 12.7, 14.2 and 14.9.
\url{https://math.uchicago.edu/~may/PEOPLE/MIKE/framed.pdf}.

\bibitem{Weber}
M. Weber. \emph{Polynomials in categories with pullbacks}.
arXiv:1106.1983. Distributivity pullbacks in Section 2.2 and polynomial
2-bicategories in Section 4.1.
\url{https://arxiv.org/abs/1106.1983}.

\bibitem{FTMH}
S. Fujii, Y. C. Tsai, Y. Montacute and I. Hasuo.
\emph{Monads and Distributive Laws in Substructural Contexts}, extended
version, 2026. Proposition 36 and Theorem 42.
\url{https://arxiv.org/html/2605.13533v1}.

\bibitem{FongZanasi}
B. Fong and F. Zanasi. \emph{Universal Constructions for (Co)Relations}.
arXiv:1710.03894v3. Assumption 2, Section 3.2, Corollary 3.4 and Section 4.4.
\url{https://arxiv.org/abs/1710.03894}.

\bibitem{RosebrughWood}
R. Rosebrugh and R. J. Wood. \emph{Distributive Laws and Factorization}.
\url{https://mta.ca/~rrosebru/articles/fac.pdf}.

\bibitem{Stepan}
M. \v{S}t\v{e}p\'an. \emph{Factorization systems and double categories}.
arXiv:2305.06714.
\url{https://arxiv.org/abs/2305.06714}.
\end{thebibliography}
\end{document}
